% =====================================================================
%  Промежуточный отчёт. Этап 1 ВКР
%
%  Тема: Методика обнаружения несанкционированных действий и аномальной
%        активности в корпоративной сети с использованием нейронных сетей
%
%  Сборка: xelatex -> biber -> xelatex -> xelatex
% =====================================================================
\documentclass[14pt, a4paper]{extarticle}

% =====================================================================
%  Подключаемые пакеты LaTeX
%  Сборка: XeLaTeX + Biber
%  Соответствие: ГОСТ 7.32-2017, ГОСТ Р 7.0.100-2018
% =====================================================================

% Кодировка и язык
\usepackage{fontspec}
\setmainfont{Times New Roman}
\setsansfont{Arial}
\setmonofont{Courier New}

\usepackage{polyglossia}
\setdefaultlanguage{russian}
\setotherlanguage{english}

% Геометрия страницы (ГОСТ 7.32-2017, п. 6.1):
% левое 30 мм, правое 15 мм, верхнее 20 мм, нижнее 20 мм
\usepackage[
    a4paper,
    left=30mm,
    right=15mm,
    top=20mm,
    bottom=20mm,
    footskip=10mm
]{geometry}

% Межстрочный интервал 1.5 (ГОСТ 7.32-2017, п. 6.3)
\usepackage{setspace}
\onehalfspacing

% Абзацный отступ 1.25 см
\usepackage{indentfirst}
\setlength{\parindent}{1.25cm}

% Математика
\usepackage{amsmath}
\usepackage{amssymb}
\usepackage{amsthm}
\usepackage{mathtools}

% Графика и таблицы
\usepackage{graphicx}
\usepackage{booktabs}
\usepackage{longtable}
\usepackage{tabularx}
\usepackage{array}
\usepackage{multirow}
\usepackage{makecell}
\usepackage{caption}
\usepackage{subcaption}
\usepackage{float}

% Цвет, ссылки, листинги
\usepackage{xcolor}
\usepackage[hidelinks]{hyperref}
\hypersetup{
    pdfencoding=auto,
    bookmarksnumbered=true,
    bookmarksopen=true,
    colorlinks=false
}
\usepackage{listings}
\lstset{
    basicstyle=\ttfamily\footnotesize,
    breaklines=true,
    showstringspaces=false,
    keywordstyle=\bfseries,
    columns=fullflexible,
    captionpos=b,
    frame=single,
    aboveskip=8pt,
    belowskip=8pt
}

% Перечисления и списки
\usepackage{enumitem}
\setlist{nosep, leftmargin=*}

% Заголовки и оглавление
\usepackage{titlesec}
\usepackage{tocloft}

% Библиография по ГОСТ Р 7.0.100-2018
\usepackage[
    backend=biber,
    style=gost-numeric,
    language=auto,
    sorting=none,
    maxbibnames=4,
    minbibnames=1
]{biblatex}

% Дополнительные служебные пакеты
\usepackage{csquotes}
\usepackage{etoolbox}
\usepackage{url}
\usepackage{microtype}
\usepackage{xurl}

% =====================================================================
%  Оформление по ГОСТ 7.32-2017
%  Структурные элементы: разделы, рисунки, таблицы, формулы, списки
% =====================================================================

% Размер основного шрифта 14 pt с межстрочным интервалом 1.5 устанавливается
% в classfile (extarticle/extreport, базовый размер 14pt). Здесь — стилевые
% настройки заголовков и нумерации.

% Заголовки по ГОСТ 7.32-2017 (п. 6.2):
%   - заголовки разделов и подразделов с абзацного отступа,
%     с прописной буквы, без точки в конце, не подчёркиваются;
%   - расстояние между заголовком и текстом — не менее одного интервала;
%   - заголовки разделов набираются жирным шрифтом, размер 14 pt.

\titleformat{\section}
    {\normalfont\fontsize{14}{17}\bfseries}
    {\thesection}{1em}{}
\titlespacing*{\section}{\parindent}{18pt}{12pt}

\titleformat{\subsection}
    {\normalfont\fontsize{14}{17}\bfseries}
    {\thesubsection}{1em}{}
\titlespacing*{\subsection}{\parindent}{14pt}{10pt}

\titleformat{\subsubsection}
    {\normalfont\fontsize{14}{17}\bfseries}
    {\thesubsubsection}{1em}{}
\titlespacing*{\subsubsection}{\parindent}{12pt}{8pt}

% Нумерация рисунков и таблиц — сквозная в пределах раздела
\renewcommand{\thefigure}{\arabic{section}.\arabic{figure}}
\renewcommand{\thetable}{\arabic{section}.\arabic{table}}
\renewcommand{\theequation}{\arabic{section}.\arabic{equation}}

% Подписи рисунков и таблиц по ГОСТ:
%   - «Рисунок N — Название» (тире), снизу, по центру;
%   - «Таблица N — Название», сверху, слева.
\captionsetup[figure]{
    labelsep=endash,
    name=Рисунок,
    justification=centering,
    singlelinecheck=true,
    font=normalsize
}
\captionsetup[table]{
    labelsep=endash,
    name=Таблица,
    justification=raggedright,
    singlelinecheck=false,
    font=normalsize
}

% Колонтитулы — простые, номер страницы по центру нижнего поля
\usepackage{fancyhdr}
\pagestyle{fancy}
\fancyhf{}
\fancyfoot[C]{\thepage}
\renewcommand{\headrulewidth}{0pt}
\renewcommand{\footrulewidth}{0pt}

% Оглавление: «СОДЕРЖАНИЕ» по центру, прописными
\renewcommand{\contentsname}{\hfill СОДЕРЖАНИЕ\hfill}

% Список литературы: «СПИСОК ИСПОЛЬЗОВАННЫХ ИСТОЧНИКОВ»
\defbibheading{bibliography}[Список использованных источников]{%
    \section*{#1}\addcontentsline{toc}{section}{#1}%
}

% =====================================================================
%  Пользовательские команды и сокращения
% =====================================================================

% Краткие ссылки
\newcommand{\figref}[1]{рисунок~\ref{#1}}
\newcommand{\Figref}[1]{Рисунок~\ref{#1}}
\newcommand{\tabref}[1]{таблица~\ref{#1}}
\newcommand{\Tabref}[1]{Таблица~\ref{#1}}
\newcommand{\eqnref}[1]{(\ref{#1})}
\newcommand{\secref}[1]{раздел~\ref{#1}}

% Английские сокращения как термины
\newcommand{\eng}[1]{\textit{#1}}

% Метки общих обозначений
\DeclareMathOperator*{\argmax}{arg\,max}
\DeclareMathOperator*{\argmin}{arg\,min}
\DeclareMathOperator{\softmax}{softmax}
\DeclareMathOperator{\sigm}{\sigma}
\DeclareMathOperator{\KL}{KL}
\DeclareMathOperator{\MMD}{MMD}
\DeclareMathOperator{\PSI}{PSI}

% Векторы и матрицы — обычным начертанием в соответствии с ГОСТ 7.32
% (без полужирного, как принято в большинстве отечественных НИР)

% Окружение для алгоритмов
\newtheoremstyle{algodef}{8pt}{8pt}{\normalfont}{}{\bfseries}{.}{0.5em}{}
\theoremstyle{algodef}
\newtheorem{algorithm}{Алгоритм}[section]
\newtheorem{definition}{Определение}[section]
\newtheorem{statement}{Утверждение}[section]


\addbibresource{bib/refs.bib}

% Гиперссылки на разделы — простые, без подсветки
\hypersetup{
    pdftitle={Промежуточный отчёт. Этап 1 ВКР},
    pdfauthor={},
    pdfsubject={Методика обнаружения несанкционированных действий и
                аномальной активности в корпоративной сети
                с использованием нейронных сетей},
    pdfkeywords={NIDS, NTA, нейронные сети, CNN-LSTM, UNSW-NB15,
                 обнаружение аномалий, ИИ-злоумышленник, drift}
}

\begin{document}

% ----------- Титульный лист (упрощённый, для отчёта) -----------
\begin{titlepage}
\centering
\vspace*{1cm}

{\large\bfseries Промежуточный отчёт по выпускной квалификационной работе}\\[0.3cm]
{\large\bfseries Этап 1. Теоретический анализ}\\[2cm]

{\Large\bfseries Тема: <<Методика обнаружения несанкционированных действий
и аномальной активности в корпоративной сети с использованием нейронных
сетей>>}\\[3cm]

\begin{flushright}
\begin{tabular}{rl}
Исполнитель: & \rule{6cm}{0.4pt} \\[0.3cm]
Научный руководитель: & \rule{6cm}{0.4pt} \\
\end{tabular}
\end{flushright}

\vfill
{\large \the\year}\\
\end{titlepage}

% ----------- Реферат (по ГОСТ 7.32-2017, п. 5.3) -----------
\section*{РЕФЕРАТ}
\addcontentsline{toc}{section}{РЕФЕРАТ}

Отчёт содержит \pageref{LastPage} страниц,
\ref{lastfig} рисунков, \ref{lasttab} таблиц,
\ref{lastref} использованных источников.

\textbf{Ключевые слова:} обнаружение вторжений, NIDS, NTA,
анализ сетевого трафика, глубокое обучение, нейронные сети, CNN-LSTM,
UNSW-NB15, CICIDS2017, ИИ-злоумышленник, дрейф распределений,
обнаружение аномалий.

\textbf{Объект исследования} --- сетевой трафик корпоративной сети,
подверженной несанкционированным действиям и аномальной активности.

\textbf{Предмет исследования} --- методика обнаружения аномальной
активности и несанкционированных действий с использованием нейронных
сетей и активного тестирования при помощи ИИ-злоумышленника.

\textbf{Цель работы} --- разработка методики обнаружения аномальной
активности и несанкционированных действий в корпоративной сети на основе
нейронных сетей, включающей анализ сетевого трафика, классификацию угроз
и моделирование атак с помощью ИИ-злоумышленника, генерирующего трафик
в реальном времени для тестирования системы.

\textbf{Содержание этапа.} В рамках Этапа 1 выполнены: систематический
анализ современных методов и технологий обнаружения аномалий и атак;
сравнительный анализ десяти публично доступных датасетов трафика и
обоснование выбора UNSW-NB15 как основного и CICIDS2017 как
вспомогательного для cross-evaluation; теоретический обзор подходов к
моделированию атак с использованием искусственного интеллекта и
обоснование архитектуры ИИ-злоумышленника; формирование функциональных
и нефункциональных требований к системе и формализация гипотезы
исследования. Полученные результаты используются как исходные данные
Этапа 2 (теоретическое сравнение архитектур и проектирование методики).

\newpage

% ----------- Перечень сокращений и обозначений -----------
\section*{ПЕРЕЧЕНЬ СОКРАЩЕНИЙ И УСЛОВНЫХ ОБОЗНАЧЕНИЙ}
\addcontentsline{toc}{section}{ПЕРЕЧЕНЬ СОКРАЩЕНИЙ И УСЛОВНЫХ ОБОЗНАЧЕНИЙ}

\begin{description}[font=\normalfont\bfseries, leftmargin=2.5em]
    \item[AE] Autoencoder, автоэнкодер
    \item[BiLSTM] Bidirectional Long Short-Term Memory
    \item[CNN] Convolutional Neural Network, свёрточная нейронная сеть
    \item[DL] Deep Learning, глубокое обучение
    \item[DPI] Deep Packet Inspection, глубокий анализ пакета
    \item[ECE] Expected Calibration Error
    \item[EPS] Events Per Second, событий в секунду
    \item[FPR] False Positive Rate
    \item[FN] False Negative
    \item[FP] False Positive
    \item[FSM] Finite-State Machine, конечный автомат
    \item[GAN] Generative Adversarial Network
    \item[GRU] Gated Recurrent Unit
    \item[IDS] Intrusion Detection System
    \item[IPFIX] IP Flow Information Export
    \item[KL] Kullback--Leibler divergence
    \item[LSTM] Long Short-Term Memory
    \item[MCC] Matthews Correlation Coefficient
    \item[ML] Machine Learning
    \item[MLP] Multilayer Perceptron
    \item[MMD] Maximum Mean Discrepancy
    \item[NIDS] Network-based Intrusion Detection System
    \item[NTA] Network Traffic Analysis
    \item[OCSVM] One-Class Support Vector Machine
    \item[PSI] Population Stability Index
    \item[PR-AUC] Area Under the Precision--Recall Curve
    \item[ReLU] Rectified Linear Unit
    \item[RF] Random Forest
    \item[RNN] Recurrent Neural Network
    \item[ROC-AUC] Area Under the Receiver Operating Characteristic Curve
    \item[SOC] Security Operations Center
    \item[SOAR] Security Orchestration, Automation and Response
    \item[SOM] Self-Organizing Map
    \item[SVM] Support Vector Machine
    \item[TCN] Temporal Convolutional Network
    \item[TN] True Negative
    \item[TP] True Positive
    \item[TPR] True Positive Rate, Recall
    \item[TTD] Time-To-Detect
    \item[VAE] Variational Autoencoder
    \item[XGBoost] eXtreme Gradient Boosting
\end{description}

\newpage

% ----------- Содержание -----------
\tableofcontents
\newpage

% ----------- Введение -----------
\section*{ВВЕДЕНИЕ}
\addcontentsline{toc}{section}{ВВЕДЕНИЕ}

\textbf{Актуальность.} Защита корпоративных сетей от несанкционированных
действий и аномальной активности является устойчивым направлением
прикладных исследований в области информационной безопасности. Рост
сложности атак, доминирование шифрованного трафика и динамика
распределений в продуктивной среде делают сигнатурные системы
недостаточными и переводят центр тяжести разработки в сторону
нейросетевых аномальных детекторов. При этом актуальной остаётся
проблема воспроизводимой и формализованной оценки таких детекторов в
условиях, приближенных к реальной эксплуатации, в том числе при дрейфе
распределений.

\textbf{Цель и задачи Этапа 1.} Целью этапа является формирование
теоретической базы исследования и обоснование принципиальных
проектных решений: выбора датасета, набора моделей-кандидатов,
архитектуры ИИ-злоумышленника, формальных требований и гипотезы.

Задачи этапа включают:
\begin{enumerate}
    \item Систематизацию методов и технологий обнаружения аномалий и
    атак, включая статистические, классические ML и DL-подходы.
    \item Сравнительный анализ публично доступных датасетов NIDS и
    обоснование выбора основного датасета.
    \item Обзор подходов к моделированию атак с применением ИИ и
    обоснование архитектуры ИИ-злоумышленника.
    \item Формирование функциональных и нефункциональных требований к
    системе и формализацию гипотезы исследования.
\end{enumerate}

\textbf{Методы исследования.} Систематический анализ литературы;
сопоставление по формализованным критериям; теоретический анализ
математических моделей; экспертная оценка ограничений рассматриваемых
подходов.

\textbf{Структура отчёта.} Раздел~1 содержит обзор IDS/NIDS, методов
обнаружения аномалий, нейросетевых архитектур и метрик качества.
Раздел~2 посвящён сравнительному анализу датасетов и обоснованию выбора
UNSW-NB15. В разделе~3 рассмотрены подходы к моделированию атак и
обоснована архитектура ИИ-злоумышленника. В разделе~4 формализованы
требования к системе и гипотеза исследования.

\newpage

% ----------- Основная часть -----------
% =====================================================================
%  Stage 1, раздел 1.
%  Обзор IDS/IPS/NIDS, методов обнаружения аномалий,
%  нейросетевых подходов и метрик качества.
% =====================================================================

\section{Корпоративная сеть как объект мониторинга и обнаружения несанкционированных действий}

\subsection{Понятие корпоративной сети и объекты мониторинга}

В рамках настоящего исследования под корпоративной сетью понимается совокупность
взаимосвязанных программно-аппаратных средств, обеспечивающих обмен данными между
рабочими станциями, серверами, сетевым оборудованием и внешними сегментами (Internet,
филиалы, облачные провайдеры). С точки зрения мониторинга безопасности
корпоративная сеть характеризуется тремя признаками, существенными для построения
системы обнаружения вторжений:

\begin{enumerate}
    \item \textbf{Высокой неоднородностью трафика.} На периметре и в сегментах LAN
    одновременно сосуществуют десятки протоколов прикладного уровня, что приводит к
    мультимодальному распределению признаков сетевых соединений и затрудняет
    применение единых пороговых правил~\cite{khraisat2019survey, ahmad2021nids}.
    \item \textbf{Доминированием легитимного трафика.} Доля атакующих соединений в
    реальной эксплуатации, как правило, не превышает долей процента; это создаёт
    ярко выраженный дисбаланс классов и приводит к необходимости использовать
    специализированные методы обучения и оценки~\cite{sommer2010outside,
    chawla2002smote, lin2017focal}.
    \item \textbf{Динамичностью характеристик.} Состав сетевых сервисов, схемы
    маршрутизации, профили активности пользователей и спектр атак изменяются во
    времени, что приводит к явлению дрейфа распределений (\eng{concept drift,
    covariate shift})~\cite{gama2014drift, andresini2021insomnia}.
\end{enumerate}

В качестве объектов мониторинга в работе рассматриваются: сетевые соединения
(\eng{flows}) на уровне L3/L4, сессии прикладного уровня и временные
последовательности соединений в скользящих окнах. Это согласуется с подходом,
принятым в большинстве современных систем сетевого анализа трафика
(\eng{Network Traffic Analysis, NTA}) и в открытых источниках
данных~\cite{moustafa2015unsw, sharafaldin2018toward}.

\subsection{Модели угроз и таксономия атак}

При проектировании NIDS целесообразно опираться на формализованные таксономии угроз.
Наиболее распространёнными являются следующие модели.

\begin{itemize}
    \item \textbf{MITRE ATT\&CK Enterprise} --- матрица тактик и техник атак,
    систематизирующая действия атакующего на стадиях разведки, первоначального
    доступа, закрепления, эскалации привилегий, латерального перемещения,
    сбора и эксфильтрации данных~\cite{mitreattack}.
    \item \textbf{Таксономия датасета UNSW-NB15}, выделяющая девять семейств
    атакующего трафика: \emph{Fuzzers, Analysis, Backdoors, DoS, Exploits, Generic,
    Reconnaissance, Shellcode, Worms}~\cite{moustafa2015unsw}.
    \item \textbf{Таксономия датасета CICIDS2017}, содержащая категории
    \emph{Brute Force, Heartbleed, Botnet, DoS, DDoS, Web Attack, Infiltration,
    Port Scan}~\cite{sharafaldin2018toward}.
\end{itemize}

В диссертационной работе в качестве рабочей таксономии принята классификация
UNSW-NB15, поскольку именно она используется для разметки основного обучающего
корпуса. Расширение на дополнительные классы атак (\textit{Brute Force, Web
Attack, Botnet}) выполняется на этапе кросс-валидации с CICIDS2017.

\subsection{Виды несанкционированной активности и их сетевые проявления}

С позиции NIDS все рассматриваемые типы вторжений интерпретируются через
проявления в признаках сетевых соединений. Сводная классификация
сетевых проявлений приведена в \tabref{tab:traffic-anomalies}.

\begin{table}[H]
\centering
\caption{Сетевые проявления типовых атак в признаках flow-уровня}
\label{tab:traffic-anomalies}
\renewcommand{\arraystretch}{1.2}
\begin{tabularx}{\linewidth}{@{}lXX@{}}
\toprule
\textbf{Класс атаки} & \textbf{Характеристики во flow-признаках} & \textbf{Признаки временной структуры} \\
\midrule
DoS / DDoS &
Резкий рост числа соединений к одному адресату; короткие потоки; рост доли
неуспешных TCP-handshake; всплески объёма пакетов в единицу времени &
Сжатие межпакетных интервалов, синхронизированные всплески входящего трафика \\
\addlinespace
Reconnaissance, Port Scan &
Большое число потоков к одному источнику с разными destination-портами;
короткие соединения; преобладание SYN/ICMP &
Регулярные интервалы между потоками, монотонный обход диапазона портов \\
\addlinespace
Brute Force &
Повторные подключения к одному порту (22, 3389, 80); ошибки авторизации в
прикладной части; малое число байт за сессию &
Высокая периодичность в окне порядка единиц секунд \\
\addlinespace
Exploits, Backdoors &
Аномальные размеры payload; нестандартные сочетания TCP-флагов;
несвойственное соотношение в обе стороны &
Резкие переходы между фазой сканирования и фазой эксплуатации \\
\addlinespace
Data Exfiltration &
Длинные потоки с высоким \(b_{out}\) и низким \(b_{in}\); устойчивая
скорость передачи; нестандартные протоколы &
Долговременные сессии с малой вариативностью скорости \\
\addlinespace
Worm, Botnet, Lateral movement &
Большое число направленных потоков от одного источника; повторение шаблонных
сессий; обращения к C\&C-инфраструктуре &
Циклическая активность с фиксированным периодом, синхронные всплески \\
\bottomrule
\end{tabularx}
\end{table}

Из \tabref{tab:traffic-anomalies} видно, что значительная часть атак проявляется
не в отдельных признаках, а в их совместном распределении и временной структуре.
Это аргументирует использование моделей, способных одновременно учитывать
локальные паттерны (свёрточные слои) и временные зависимости (рекуррентные слои
или механизмы внимания)~\cite{kim2016lstm, yin2017dlrnn, vinayakumar2019dl}.

\section{Системы обнаружения вторжений: классификация и принципы построения}

\subsection{Классификация IDS по принципу детекции}

В соответствии с устоявшейся в литературе типологией~\cite{khraisat2019survey,
ahmad2021nids, ferrag2020dlcyber} системы обнаружения вторжений делятся на три
крупных класса.

\begin{enumerate}
    \item \textbf{Сигнатурные (\eng{signature-based, misuse detection}).}
    Принимают решение по совпадению наблюдаемого поведения с заранее заданным
    шаблоном (сигнатурой). Преимущества: высокая точность по известным угрозам,
    интерпретируемость, низкая частота ложных срабатываний. Ограничения:
    неприменимость к ранее не встречавшимся атакам (\eng{zero-day}),
    необходимость поддержки баз сигнатур, уязвимость к обфускации трафика.
    Промышленные представители: Snort, Suricata, Zeek (в части сигнатурных
    модулей).
    \item \textbf{Аномальные (\eng{anomaly-based}).} Оценивают отклонение
    наблюдаемого поведения от модели нормы и сигнализируют при превышении
    порога. Преимущества: способность детектировать ранее неизвестные атаки.
    Ограничения: повышенная частота ложных срабатываний, чувствительность к
    дрейфу нормального профиля и к качеству обучающей выборки. Подкласс
    \eng{specification-based} использует формальные модели нормы.
    \item \textbf{Гибридные.} Комбинируют сигнатурную и аномальную ветви,
    распределяя решения по уровню риска и стоимости ошибок: сигнатурная ветвь
    подтверждает известные классы, аномальная --- расширяет покрытие на
    неизвестные сценарии.
\end{enumerate}

В настоящем исследовании рассматривается аномальный класс с акцентом на
нейросетевые модели; гибридная конфигурация выступает рекомендуемой формой
интеграции в SOC, описанной в \secref{subsec:soc-integration}.

\subsection{Packet-based и flow-based анализ}

С точки зрения исходного представления данных различают два уровня анализа.

\textbf{Packet-based (DPI, Deep Packet Inspection).} Объект анализа --- отдельный
сетевой пакет вместе с его полезной нагрузкой. Преимущества: максимальный объём
семантической информации, возможность распознавать атаки на прикладном уровне.
Ограничения: высокая ресурсоёмкость, сложность шифрованного трафика, ограничения
по приватности, затрудняющие использование в корпоративной среде.

\textbf{Flow-based.} Объект анализа --- сетевое соединение (flow), описываемое
агрегированными характеристиками: длительность, количество пакетов и байтов в
обоих направлениях, флаги TCP, межпакетные интервалы, протокол и сервис.
Стандартными протоколами экспорта flow-записей являются Cisco
NetFlow~\cite{rfc3954} и IETF IPFIX~\cite{rfc7011}. Преимущества подхода:
ресурсная эффективность, устойчивость к шифрованию (анализ метаданных), удобство
агрегирования в скользящие окна.

В работе используется flow-based подход. Обоснование: датасеты UNSW-NB15 и
CICIDS2017 представлены именно в этом формате; flow-уровень напрямую
соответствует выходу типовых сборщиков NetFlow/IPFIX, что обеспечивает
переносимость методики в реальные инфраструктуры~\cite{moustafa2015unsw,
sharafaldin2018toward, gharib2016evalframework}.

\subsection{Архитектура SOC и роль NIDS/NTA в общем pipeline алертирования}
\label{subsec:soc-integration}

Современный SOC реализуется как многоуровневый конвейер обработки телеметрии:
сборщики (\eng{network sensors, EDR, log forwarders}) $\rightarrow$ потоковые
очереди (\eng{Kafka, NATS}) $\rightarrow$ нормализаторы $\rightarrow$ хранилища
(SIEM, data lake) $\rightarrow$ аналитические подсистемы (NIDS, UEBA, корреляторы)
$\rightarrow$ модули реагирования (SOAR). Сетевой детектор аномалий выступает
аналитической подсистемой и обязан удовлетворять трём ограничениям:

\begin{itemize}
    \item ограниченное время отклика (\eng{time-to-alert}), достаточное для
    реактивного блокирования;
    \item низкая частота ложных срабатываний на единицу времени
    (\eng{false positives per time}), поскольку каждое срабатывание формирует
    задание для аналитика SOC;
    \item возможность приоритизации алертов, чтобы аналитик в первую очередь
    рассматривал инциденты с высоким влиянием.
\end{itemize}

Указанные ограничения определяют требования к метрикам, рассматриваемым далее в
\secref{sec:metrics}, и к схеме формирования алертов прототипа, описываемой на
этапе 3 настоящей работы.

\section{Методы обнаружения аномалий}

В этом разделе систематизирована теоретическая база основных групп методов,
применяемых для детекции аномалий, и описаны их математические основы.

\subsection{Постановка задачи обнаружения аномалий}

Пусть \(X \subseteq \mathbb{R}^d\) --- пространство признаков сетевых соединений
(или окон), \(p_{0}(x)\) --- плотность распределения нормальных соединений.
Задача состоит в построении функции скоринга \(s\colon X \to \mathbb{R}\) и
порога \(\tau \in \mathbb{R}\) такими, что
\begin{equation}
\hat{y}(x) = \mathbb{I}\bigl[s(x) > \tau\bigr],
\quad x \in X,
\label{eq:detector}
\end{equation}
где \(\hat{y}(x) = 1\) интерпретируется как ``аномалия / атака''.
Обучение проводится в одном из трёх режимов:

\begin{itemize}
    \item \textbf{Supervised}: дана выборка \(\{(x_i, y_i)\}_{i=1}^{n}\) с
    метками классов \(y_i \in \{0,1\}\); скоринг строится через минимизацию
    эмпирического риска.
    \item \textbf{Semi-supervised (one-class)}: даны только нормальные
    наблюдения; скоринг строится как мера отклонения от модели нормы.
    \item \textbf{Unsupervised}: выборка немаркирована; скоринг получается из
    статистических, плотностных или структурных свойств данных.
\end{itemize}

Каждый из режимов применим в условиях NIDS и обладает собственными
ограничениями~\cite{khraisat2019survey, sommer2010outside}. В дальнейшем
исследовании используется преимущественно supervised-режим (на UNSW-NB15) с
параллельным включением semi-supervised-моделей на базе автоэнкодеров и
one-class классификаторов в качестве сравнения.

\subsection{Статистические и энтропийные методы}

Простейший статистический подход использует параметрическое описание
распределения признаков. Например, для одномерного признака \(x\) типичной
является \(z\)-оценка
\begin{equation}
z = \frac{x - \mu}{\sigma},
\qquad
\hat{y} = \mathbb{I}\bigl[\,|z| > \tau_z\bigr],
\label{eq:zscore}
\end{equation}
где \(\mu\), \(\sigma\) --- оценки математического ожидания и
среднеквадратического отклонения. Многомерное обобщение --- расстояние
Махаланобиса
\begin{equation}
d_M(x) = \sqrt{(x-\mu)^{\top}\Sigma^{-1}(x-\mu)},
\label{eq:mahalanobis}
\end{equation}
где \(\Sigma\) --- ковариационная матрица. Подход эффективен при
гауссовости признаков и небольшой размерности, но в реальном NIDS-контексте
сталкивается с многомодальностью распределений и тяжёлыми хвостами.

Энтропийные методы оценивают энтропию распределений ключевых признаков
(распределение IP-адресов, портов, размеров пакетов) в скользящем окне,
интерпретируя резкое изменение энтропии как аномалию. Для дискретного
распределения \(P\) над алфавитом \(\mathcal{A}\) энтропия Шеннона определяется
как
\begin{equation}
H(P) = -\sum_{a \in \mathcal{A}} P(a)\log P(a).
\label{eq:entropy}
\end{equation}
Подобные индикаторы исторически использовались в подсистемах NetFlow-анализа,
но в современных условиях, как отмечается в обзорах~\cite{khraisat2019survey,
ferrag2020dlcyber}, их разрешающая способность недостаточна для распределённых
и низкочастотных атак.

\subsection{Кластеризация и плотностные методы}

\textbf{K-means} формирует разбиение на \(K\) кластеров с центроидами
\(\{c_k\}\) минимизацией внутрикластерной дисперсии:
\begin{equation}
\min_{\{c_k\}}\sum_{k=1}^{K}\sum_{x \in C_k}\|x - c_k\|^{2}.
\label{eq:kmeans}
\end{equation}
Аномалия определяется через расстояние до ближайшего центроида.

\textbf{DBSCAN} определяет аномалии как точки, не принадлежащие ни одному
плотностному кластеру. Параметры: радиус окрестности \(\varepsilon\) и
минимальное число точек \(\mathit{minPts}\). Метод устойчив к
неэллипсоидальной форме кластеров, но плохо масштабируется на большие выборки.

\textbf{LOF (Local Outlier Factor)} вычисляет относительную плотность точки
по сравнению с её \(k\)-ближайшими соседями. Точка с локальной плотностью,
существенно ниже плотности соседей, считается выбросом.

\textbf{Isolation Forest}~\cite{liu2008isolation} использует случайное
рекурсивное разбиение пространства; чем меньше глубина, на которой точка
изолируется, тем выше её аномальность. Для узла-листа с глубиной \(h(x)\)
скоринг определяется как
\begin{equation}
s(x, n) = 2^{-\frac{\mathbb{E}[h(x)]}{c(n)}},
\qquad
c(n) \approx 2\bigl(\ln(n-1) + 0.5772\bigr) - \frac{2(n-1)}{n},
\label{eq:iforest}
\end{equation}
где \(c(n)\) --- средняя длина пути в случайном бинарном дереве из \(n\)
точек. Метод отличается линейной по объёму данных сложностью и хорошо
переносится на признаковые пространства NIDS.

\subsection{One-class классификация}

\textbf{One-Class SVM}~\cite{scholkopf2001ocsvm} ищет в спрямляющем
пространстве минимальную гиперсферу, содержащую долю \(1-\nu\) обучающих
наблюдений, что эквивалентно построению разделяющей гиперплоскости с
максимальным отступом от начала координат:
\begin{equation}
\min_{w, \rho, \xi} \tfrac{1}{2}\|w\|^{2} + \tfrac{1}{\nu n}\sum_{i=1}^{n}\xi_i - \rho,
\quad \text{s.t.}\quad \langle w, \phi(x_i)\rangle \geq \rho - \xi_i,\ \xi_i \geq 0,
\label{eq:ocsvm}
\end{equation}
где \(\phi\) --- спрямляющее преобразование, индуцированное ядром.
Аномалия идентифицируется как отрицательное значение функции
\(f(x) = \langle w, \phi(x)\rangle - \rho\).

Параметр \(\nu \in (0, 1]\) задаёт верхнюю оценку доли выбросов в обучающей
выборке. Метод чувствителен к выбору ядра (радиальное гауссовское ядро
наиболее распространено) и требует тщательного масштабирования признаков.

\subsection{Нейросетевые методы детекции аномалий}

Нейросетевые модели в задачах NIDS реализуются в одном из двух режимов:
дискриминативном (бинарная или мультиклассовая классификация) и
реконструкционно-вероятностном (автоэнкодеры, вариационные автоэнкодеры).
Принципиальное отличие от классических методов состоит в способности
автоматически извлекать иерархические признаки из исходных flow-векторов и
последовательностей соединений~\cite{lecun2015dl, vinayakumar2019dl,
mirsky2018kitsune}. Подробный анализ архитектур и обоснование выбора --- в
\secref{sec:dl-architectures}.

\section{Глубинное обучение для анализа сетевого трафика}
\label{sec:dl-architectures}

\subsection{Полносвязные нейронные сети (MLP)}

Многослойный перцептрон выступает базовой моделью DL и описывается
рекурсией
\begin{equation}
h^{(l)} = \phi\!\left(W^{(l)} h^{(l-1)} + b^{(l)}\right),
\qquad l = 1, \dots, L,
\label{eq:mlp}
\end{equation}
где \(\phi\) --- нелинейность (ReLU, LeakyReLU, GELU), \(W^{(l)}\), \(b^{(l)}\)
--- обучаемые параметры. Для бинарной классификации выход формируется как
\(\hat{p}(y=1\mid x) = \sigma(W^{(L+1)} h^{(L)} + b^{(L+1)})\). Обучение
сводится к минимизации бинарной перекрёстной энтропии
\begin{equation}
\mathcal{L}_{\text{BCE}} =
-\frac{1}{n}\sum_{i=1}^{n}\Bigl[y_i \log \hat{p}_i + (1-y_i)\log(1-\hat{p}_i)\Bigr].
\label{eq:bce}
\end{equation}
В контексте NIDS MLP служит нижней границей DL-моделей: он не учитывает ни
локальной, ни временной структуры, но позволяет проверить достаточность
признакового представления.

\subsection{Свёрточные сети для последовательностей (1D-CNN)}

Одномерные свёрточные слои применяются для извлечения локальных паттернов из
последовательностей flow-признаков в скользящем окне. Операция
свёртки имеет вид
\begin{equation}
(h * w)_t = \sum_{k=0}^{K-1} w_k\, h_{t-k},
\label{eq:conv1d}
\end{equation}
где \(K\) --- размер ядра. Стек свёрточных слоёв с пуллингом формирует
иерархию рецептивных полей, позволяя обнаруживать локальные ``мотивы''
поведения, характерные для конкретных классов атак (например, пакеты SYN-флуда
в окне порядка десятков соединений). По данным
работ~\cite{vinayakumar2019dl, ferrag2020dlcyber}, 1D-CNN превосходит MLP по
полноте детекции атак типа DoS и Reconnaissance.

\subsection{Рекуррентные сети: RNN, LSTM, GRU, BiLSTM}

Рекуррентная сеть моделирует временные зависимости через скрытое состояние
\(h_t\):
\begin{equation}
h_t = \phi(W_{xh} x_t + W_{hh} h_{t-1} + b_h).
\label{eq:rnn}
\end{equation}
Простая RNN страдает от проблемы исчезающего/взрывающегося градиента, что
делает её малопригодной для длинных последовательностей.

\textbf{LSTM}~\cite{hochreiter1997lstm} вводит ячейку состояния
\(c_t\) и три гейта (входной \(i_t\), забывания \(f_t\), выходной \(o_t\)):
\begin{align}
i_t &= \sigma(W_i [h_{t-1}, x_t] + b_i),\\
f_t &= \sigma(W_f [h_{t-1}, x_t] + b_f),\\
o_t &= \sigma(W_o [h_{t-1}, x_t] + b_o),\\
\tilde{c}_t &= \tanh(W_c [h_{t-1}, x_t] + b_c),\\
c_t &= f_t \odot c_{t-1} + i_t \odot \tilde{c}_t,\\
h_t &= o_t \odot \tanh(c_t).
\label{eq:lstm}
\end{align}
Гейтовая структура обеспечивает устойчивое распространение градиента и
способность модели хранить долгосрочный контекст --- свойство, востребованное
при обнаружении медленных атак (\eng{slow brute-force}, \eng{low-and-slow
exfiltration}).

\textbf{GRU}~\cite{cho2014gru} представляет собой упрощённый вариант с двумя
гейтами (\(r_t\), \(z_t\)) и без отдельной ячейки состояния. По числу
параметров GRU экономичнее LSTM при сопоставимом качестве на коротких
последовательностях.

\textbf{BiLSTM} обучает два LSTM в прямом и обратном направлениях, что
позволяет учитывать контекст ``из будущего'' при обработке элемента в
середине окна. Применимо к режимам пост-фактум анализа; для
real-time-инференса требует обработки полного окна.

\subsection{Гибридные модели CNN--LSTM}

Комбинация одномерной свёртки с рекуррентным слоем сочетает преимущества обеих
архитектур: CNN извлекает локальные паттерны (\eng{spatial features}) в
скользящем окне flow-признаков, LSTM моделирует временные зависимости между
получаемыми CNN представлениями~\cite{vinayakumar2019dl, kim2016lstm}. Схема
обработки имеет вид
\[
x_t \;\xrightarrow{\,\text{Conv1D}\,}\; z_t \;\xrightarrow{\,\text{LSTM}\,}\; h_t
\;\xrightarrow{\,\text{FC}\,}\; \hat{p}_t.
\]
Такая комбинация согласуется с природой сетевого трафика, в котором аномалия
может проявляться одновременно как ``локальный шаблон'' (последовательность
соседних соединений) и ``временной тренд'' (продолжительная активность). По
этой причине CNN--LSTM выбран в настоящей работе в качестве основной
\eng{proposed}-архитектуры (детальное обоснование --- в главе~2).

\subsection{TCN и Transformer-кодировщики}

\textbf{TCN (Temporal Convolutional Network)}~\cite{bai2018tcn} использует
причинные дилатированные свёртки и полностью обходится без рекуррентных
связей; рецептивное поле растёт экспоненциально по числу слоёв,
\(R = 1 + (K-1)\sum_{l=0}^{L-1} 2^l\). Преимущества: высокая параллелизуемость,
устойчивость градиента. Недостаток: фиксированное рецептивное поле, что в
NIDS может приводить к занижению длинных зависимостей.

\textbf{Transformer-encoder}~\cite{vaswani2017attention} основан на
механизме само-внимания
\begin{equation}
\operatorname{Attention}(Q, K, V) =
\operatorname{softmax}\!\left(\frac{QK^{\top}}{\sqrt{d_k}}\right) V,
\label{eq:attention}
\end{equation}
где \(Q\), \(K\), \(V\) --- проекции входной последовательности.
Само-внимание обеспечивает прямой доступ ко всему окну и эффективное
моделирование длинных зависимостей. Применение к NIDS требует учёта позиции
(\eng{positional encoding}) и существенно зависит от объёма обучающих данных:
по данным обзоров~\cite{ferrag2020dlcyber, ahmad2021nids}, для UNSW-NB15
TT-модели показывают результат, сопоставимый с CNN--LSTM, но с большей
дисперсией по seed.

\subsection{Автоэнкодеры и вариационные автоэнкодеры}

Автоэнкодер обучается восстанавливать вход \(x\) через узкое горлышко
\(z = E(x)\):
\begin{equation}
\mathcal{L}_{\text{AE}} = \frac{1}{n}\sum_{i=1}^{n} \|x_i - D(E(x_i))\|^{2}.
\label{eq:ae}
\end{equation}
Аномалия идентифицируется как наблюдение с ошибкой реконструкции, превышающей
порог. На таком принципе построен Kitsune~\cite{mirsky2018kitsune} ---
эталонная online-модель ансамбля автоэнкодеров для NIDS.

Вариационный автоэнкодер (VAE)~\cite{kingma2014vae} формализует кодирование
как вероятностный процесс
\begin{equation}
\mathcal{L}_{\text{VAE}} =
\mathbb{E}_{q(z|x)}[\log p(x|z)] - \KL\!\left(q(z|x)\,\|\,p(z)\right),
\label{eq:vae}
\end{equation}
где \(p(z)\) --- априорное распределение скрытого пространства,
\(q(z|x)\) --- вариационный апостериор. VAE позволяет получать вероятностный
скоринг и оценивать неопределённость предсказаний; в NIDS-применениях
демонстрирует устойчивость к шумам в обучающей выборке~\cite{lopez2017cvae}.

\subsection{Сводное теоретическое сравнение архитектур}

Сводное сопоставление, основанное на теоретических свойствах и
экспериментальных данных, опубликованных в обзорных
работах~\cite{khraisat2019survey, ahmad2021nids, liu2019survey,
ferrag2020dlcyber, vinayakumar2019dl}, представлено в \tabref{tab:dl-arch-compare}.
Численные диапазоны метрик не приводятся, поскольку зависят от датасета,
схемы оценки и не должны вырываться из контекста; раздел носит качественный
характер. Количественное сравнение выполняется в главе~2 на основании
повторяемых экспериментов.

\begin{table}[H]
\centering
\caption{Качественное сравнение нейросетевых архитектур для NIDS}
\label{tab:dl-arch-compare}
\renewcommand{\arraystretch}{1.2}
\begin{tabularx}{\linewidth}{@{}lXXXX@{}}
\toprule
\textbf{Модель} & \textbf{Локальные паттерны} & \textbf{Временные зависимости} &
\textbf{Сложность обучения} & \textbf{Inference latency} \\
\midrule
MLP        & нет          & нет        & низкая  & низкая \\
1D-CNN     & да           & ограничено & средняя & низкая \\
LSTM       & нет          & да         & средняя & средняя \\
GRU        & нет          & да         & средняя & низкая \\
BiLSTM     & нет          & двунаправ. & высокая & средняя \\
CNN--LSTM  & да           & да         & высокая & средняя \\
TCN        & да           & фиксиров.  & средняя & низкая \\
Transformer& через attention & да      & высокая & средняя \\
AE/VAE     & косвенно     & нет        & средняя & низкая \\
\bottomrule
\end{tabularx}
\end{table}

Из таблицы следует, что CNN--LSTM покрывает оба измерения структуры flow-данных
(локальное и временное) при разумной вычислительной стоимости. Дополнительные
архитектуры (BiLSTM, Transformer-encoder, TCN) рассматриваются как
конкурентные, и их прямое сопоставление приводится в главе~2.

\section{Метрики качества и протоколы оценки}
\label{sec:metrics}

\subsection{Базовые классификационные метрики}

Для бинарной задачи введём матрицу ошибок (\eng{confusion matrix}) с величинами
\(TP\), \(TN\), \(FP\), \(FN\). Тогда
\begin{align}
\text{Precision} &= \frac{TP}{TP + FP}, \label{eq:precision}\\
\text{Recall (TPR)} &= \frac{TP}{TP + FN}, \label{eq:recall}\\
\text{F1} &= \frac{2 \cdot \text{Precision} \cdot \text{Recall}}
                 {\text{Precision} + \text{Recall}},\label{eq:f1}\\
\text{FPR} &= \frac{FP}{FP + TN}.\label{eq:fpr}
\end{align}

\textbf{Precision} измеряет долю истинных атак среди вызвавших алерт; имеет
прямую интерпретацию в контексте SOC --- доля ``полезных'' алертов.
\textbf{Recall} (он же TPR) измеряет долю обнаруженных атак среди всех
реальных. \textbf{F1} --- гармоническое среднее, удобное при дисбалансе
классов. \textbf{FPR} --- долю неверно сработавших правил среди всех
нормальных наблюдений.

В условиях NIDS дополнительно используется метрика
\(\text{FP-per-time}\) --- ожидаемое число ложных срабатываний в единицу
времени, более показательное при больших объёмах нормального трафика.

Коэффициент \textbf{MCC} (Matthews correlation coefficient) учитывает все
четыре элемента матрицы ошибок и рекомендуется при сильном дисбалансе:
\begin{equation}
\text{MCC} =
\frac{TP \cdot TN - FP \cdot FN}
{\sqrt{(TP+FP)(TP+FN)(TN+FP)(TN+FN)}}.
\label{eq:mcc}
\end{equation}

\subsection{ROC-AUC, PR-AUC и кривые качества}

Поскольку детектор формирует непрерывный скоринг \(s(x)\), его качество
характеризуется не только в одной рабочей точке, но и интегрально.

\textbf{ROC-AUC} --- площадь под кривой, отображающей TPR от FPR при
изменении порога \(\tau\). Для несбалансированных задач даёт оптимистическую
оценку~\cite{khraisat2019survey, ahmad2021nids}.

\textbf{PR-AUC} --- площадь под кривой Precision--Recall. Она более
чувствительна к редкому классу и поэтому считается предпочтительной для
NIDS-сценариев с малой долей атак~\cite{liu2019survey}. В работе
PR-AUC используется как ключевая интегральная метрика качества модели.

\subsection{Эксплуатационные метрики}

К эксплуатационным метрикам относятся:
\begin{itemize}
    \item \textbf{Time-to-detect (TTD)} --- время с момента появления атаки до
    срабатывания алерта, выраженное в единицах окон или в физическом времени;
    \item \textbf{Inference latency} --- время выработки скоринга на одно
    окно;
    \item \textbf{Throughput} --- число обработанных событий в секунду
    (EPS);
    \item \textbf{Resource footprint} --- потребление CPU/GPU/RAM.
\end{itemize}
Эти величины определяют применимость модели в SOC и подбираются под целевую
эксплуатационную нагрузку.

\subsection{Калибровка и приоритизация}

Большинство нейросетевых моделей дают плохо калиброванные вероятности
\(\hat{p}\)~\cite{guo2017calibration}: уверенность модели систематически
смещена. Для использования \(\hat{p}\) в формировании приоритетов алертов
применяется температурное масштабирование
\begin{equation}
\hat{p}_{T}(y=1\mid x) =
\sigma\!\left(\frac{z(x)}{T}\right),
\label{eq:temp-scaling}
\end{equation}
где \(z(x)\) --- логит, \(T\) --- параметр, подбираемый на валидации.
Для оценки качества калибровки используется \textbf{ECE (Expected Calibration
Error)}~\cite{guo2017calibration}.

\subsection{Метрики дрейфа распределений}

Для контроля стабильности входного распределения используются:

\begin{itemize}
    \item \textbf{PSI (Population Stability Index)}: при разбиении значений
    признака на бины \(b\)
    \begin{equation}
    \PSI = \sum_{b}\,(p_b - q_b)\ln\frac{p_b}{q_b},
    \label{eq:psi}
    \end{equation}
    где \(p_b\), \(q_b\) --- доли наблюдений в бине \(b\) для эталонного и
    нового окон;
    \item \textbf{KL-расхождение}
    \(\KL(P\|Q) = \sum_b p_b \ln \tfrac{p_b}{q_b}\);
    \item \textbf{MMD (Maximum Mean Discrepancy)} --- мера различия между
    распределениями в воспроизводящем гильбертовом пространстве,
    \begin{equation}
    \MMD^{2}(P, Q) =
    \mathbb{E}_{x,x'\sim P}[k(x,x')] - 2\mathbb{E}_{x\sim P, y\sim Q}[k(x,y)]
    + \mathbb{E}_{y,y'\sim Q}[k(y,y')].
    \label{eq:mmd}
    \end{equation}
\end{itemize}

В рамках экспериментов главы~4 указанные метрики применяются для оценки
эффекта drift-инжектора и стабильности модели на потоковом сценарии.

\section{Ограничения существующих решений и постановка задачи исследования}

Системный анализ публикаций последних лет выявил следующие нерешённые
проблемы, релевантные для настоящей работы.

\begin{enumerate}
    \item Большинство экспериментов в литературе опирается на классические
    датасеты (KDD99, NSL-KDD), не отражающие современные классы атак, что
    приводит к завышенным оценкам качества~\cite{tavallaee2009nslkdd,
    moustafa2016eval, sommer2010outside, gharib2016evalframework}.
    \item Сравнение моделей часто выполняется без явного контроля дрейфа
    распределений и без оценки устойчивости при смещении тестовой выборки
    относительно обучающей~\cite{pendlebury2019tesseract,
    andresini2021insomnia}.
    \item Отсутствует общедоступный воспроизводимый стенд, в рамках которого
    модель детекции и активный (генеративный) источник атак работают в режиме
    реального времени и допускают параметризованное переключение режимов
    ``норма/атака/дрейф''~\cite{ring2019flowgan, apruzzese2020adversarial}.
    \item Большая часть исследований сосредоточена на оптимизации бинарной
    классификации и не учитывает требований SOC к приоритизации алертов и
    стоимости ложных срабатываний на единицу времени~\cite{sommer2010outside}.
\end{enumerate}

Перечисленные пробелы определяют целевые задачи диссертационной работы:
обоснованный выбор архитектуры на основе теоретического и экспериментального
сравнения; построение методики обнаружения, явно учитывающей дрейф
распределений; реализация ИИ-злоумышленника в форме воспроизводимого
шаблонного агента с конечно-автоматной политикой режимов.

\section{Промежуточные выводы по разделу}

\begin{enumerate}
    \item Корпоративная сеть характеризуется неоднородностью трафика,
    дисбалансом легитимного и атакующего поведения и динамичностью
    распределений. Это диктует необходимость использовать flow-уровень
    представления данных и модели, способные учитывать совместную и временную
    структуру признаков.
    \item Сигнатурный, аномальный и гибридный подходы NIDS дополняют, а не
    заменяют друг друга. Аномальный подход на основе нейросетей выбран в
    работе по причине способности обнаруживать ранее неизвестные сценарии
    при принципиальной возможности интеграции с сигнатурной ветвью на
    уровне SOC.
    \item Среди методов обнаружения аномалий нейросетевые модели обладают
    наибольшим потенциалом за счёт автоматического выделения иерархических
    признаков. Среди архитектур наиболее адекватной для flow-последовательностей
    является CNN--LSTM, что согласуется с теоретическим анализом и обзорной
    литературой и принимается в качестве основной (\eng{proposed}) модели.
    \item Корректная оценка качества требует одновременного использования
    интегральных (PR-AUC, MCC), пороговых (Precision, Recall, F1) и
    эксплуатационных (FP-per-time, time-to-detect, latency) метрик; для
    оценки устойчивости введены метрики дрейфа PSI, KL и MMD.
    \item Систематизация ограничений существующих решений подтверждает
    актуальность задач, заявленных в настоящей работе: построение методики,
    объединяющей нейросетевую детекцию с воспроизводимым активным
    тестированием и явной оценкой устойчивости к дрейфу.
\end{enumerate}

% =====================================================================
%  Stage 1, раздел 2.
%  Сравнительный анализ публично доступных датасетов NIDS
%  и обоснование выбора UNSW-NB15.
% =====================================================================

\section{Сравнительный анализ публично доступных датасетов NIDS}

\subsection{Критерии оценки качества датасетов}

В литературе по сетевой безопасности систематизированы критерии, по которым
оценивается пригодность датасета для задач обучения и тестирования
NIDS~\cite{gharib2016evalframework, sharafaldin2018toward, ferrag2020dlcyber}.
В настоящей работе принят следующий набор критериев, объединяющий упомянутые
источники.

\begin{enumerate}[label=K\arabic*., leftmargin=2.5em]
    \item \textbf{Актуальность.} Соответствие представленных типов атак
    современной модели угроз: наличие классов \emph{Exploits, Backdoors,
    Reconnaissance, Botnet, DDoS}; временной горизонт записи трафика;
    репрезентативность прикладных сервисов.
    \item \textbf{Репрезентативность нормального трафика.} Наличие
    реалистичного фонового поведения с разнородными прикладными протоколами.
    \item \textbf{Качество разметки.} Соответствие меток реальной природе
    трафика; согласованность таксономии; отсутствие систематических ошибок
    маркировки.
    \item \textbf{Дисбаланс классов.} Доля атакующих наблюдений и её
    распределение по типам атак; чрезмерный или искусственно сбалансированный
    дисбаланс одинаково ухудшает обобщающую способность.
    \item \textbf{Пригодность для DL.} Объём, плотность временных меток,
    наличие признаковых представлений, допускающих формирование скользящих
    окон и последовательностей.
    \item \textbf{Пригодность для flow-анализа.} Наличие готовых flow-признаков
    или возможность их извлечения; согласованность со стандартами NetFlow/IPFIX.
    \item \textbf{Воспроизводимость.} Открытость доступа, фиксированность
    официальных train/test-разделений, наличие публикаций с описанием схемы
    сбора.
    \item \textbf{Ограничения и известные дефекты.} Зафиксированные в
    литературе аномалии конкретного датасета, способные исказить результаты
    обучения и оценки.
\end{enumerate}

В дальнейшем сравнении каждый датасет описывается с точки зрения этих
критериев, после чего проводится сводная оценка.

\subsection{KDDCup99}

KDDCup99~\cite{tavallaee2009nslkdd} --- историческое расширение трафика
DARPA 1998. Содержит 41 признак (включая host-based и time-based агрегации)
и четыре крупных класса атак (\emph{DoS, Probe, R2L, U2R}).

\begin{description}[font=\normalfont\itshape, leftmargin=1.5em]
    \item[Сильные стороны:] длительная история использования; сравнимость с
    большим числом ранних работ; малый объём, удобный для отладки.
    \item[Ограничения:] устаревшая модель угроз (атаки 1998 г.); существенное
    дублирование записей в train- и test-наборах; искусственная природа
    нормального трафика, не соответствующая корпоративным сценариям;
    статистические искажения, документированные в~\cite{tavallaee2009nslkdd,
    sommer2010outside}.
\end{description}

Согласно установившейся в литературе оценке~\cite{moustafa2016eval,
ferrag2020dlcyber}, использование KDDCup99 приводит к завышенным метрикам
(близким к насыщению) и не отражает современных классов атак, что делает
датасет малопригодным для обучения промышленно ориентированной системы.

\subsection{NSL-KDD}

NSL-KDD~\cite{tavallaee2009nslkdd} --- очищенная версия KDDCup99, полученная
устранением дубликатов и балансировкой долей классов в train/test-наборах.

\begin{description}[font=\normalfont\itshape, leftmargin=1.5em]
    \item[Сильные стороны:] устранены наиболее очевидные искажения KDDCup99;
    компактный объём; удобство для классических ML-методов и для
    воспроизводимых учебных экспериментов.
    \item[Ограничения:] таксономия и природа трафика идентичны KDDCup99 ---
    атаки и сервисы 1998 г.; не содержит современных классов атак (Exploits,
    Backdoors, Botnet); ограниченная пригодность для глубинных моделей из-за
    относительно небольшого объёма (\(\approx 125\) тыс. записей в KDDTrain+).
\end{description}

NSL-KDD принят в качестве вспомогательного датасета только для исторического
сравнения с литературой; в качестве основного обучающего корпуса он не
рассматривается.

\subsection{UNSW-NB15}

UNSW-NB15~\cite{moustafa2015unsw, moustafa2016eval, unswportal} был
сформирован на стенде \emph{IXIA PerfectStorm} в Cyber Range Lab Australian
Centre for Cyber Security. Содержит \(\approx 2{,}54\,\text{млн}\) записей
сетевых соединений, описанных 49 признаками, и девять семейств атакующего
трафика: \emph{Fuzzers, Analysis, Backdoors, DoS, Exploits, Generic,
Reconnaissance, Shellcode, Worms}~\cite{moustafa2015unsw}.

\begin{description}[font=\normalfont\itshape, leftmargin=1.5em]
    \item[Сильные стороны:]
    \begin{itemize}
        \item Современная (на момент публикации) таксономия атак, включающая
        классы \emph{Exploits} и \emph{Backdoors}, которых нет в KDD-семействе.
        \item Реалистичный нормальный трафик, генерируемый в условиях
        полнопроточного эмулятора прикладных сервисов.
        \item Признаки рассчитаны как \emph{flow-features} (длительность,
        счётчики байтов и пакетов, флаги TCP, межпакетные интервалы),
        совместимые со стандартом IPFIX~\cite{rfc7011}.
        \item Зафиксированы официальные разделения \emph{UNSW\_NB15\_training-set}
        и \emph{UNSW\_NB15\_testing-set}, упрощающие сравнение между
        работами~\cite{moustafa2016eval}.
        \item Богатый набор связанных публикаций
        автора~\cite{moustafa2016eval, koroniotis2019botiot,
        moustafa2021toniot}, обеспечивающий преемственность таксономии.
    \end{itemize}
    \item[Ограничения:]
    \begin{itemize}
        \item Естественный дисбаланс классов с доминированием класса
        \emph{Generic} в составе атакующих записей.
        \item Часть признаков (\emph{ct\_*}) рассчитывается как агрегаты по
        большому контекстному окну, что ограничивает их интерпретируемость в
        реальном времени.
        \item Часть мультиклассовых меток требует дополнительной валидации;
        фигурируют единичные противоречия в распределении \emph{Analysis} ---
        это отмечено в обзорных
        публикациях~\cite{moustafa2016eval}.
    \end{itemize}
\end{description}

UNSW-NB15 удовлетворяет всем восьми критериям K1--K8 в наибольшей степени и
выбран в работе как основной обучающий и тестирующий корпус
(\secref{subsec:unsw-justification}).

\subsection{CICIDS2017}

CICIDS2017~\cite{sharafaldin2018toward} был сформирован в Canadian Institute
for Cybersecurity (CIC) на стенде с эмуляцией поведения 25 пользователей и
содержит трафик за пять рабочих дней. Атаки представлены классами
\emph{Brute Force, Heartbleed, Botnet, DoS, DDoS, Web Attack, Infiltration,
Port Scan}.

\begin{description}[font=\normalfont\itshape, leftmargin=1.5em]
    \item[Сильные стороны:] современная таксономия с акцентом на прикладные
    атаки (Brute Force, Web Attack); подробная разметка по дням; flow-признаки
    извлекаются официальной утилитой CICFlowMeter и совместимы с IPFIX.
    \item[Ограничения:] в литературе документированы дефекты разметки и
    дубликаты в исходных
    CSV-файлах~\cite{gharib2016evalframework, ferrag2020dlcyber};
    значительная доля атакующего трафика для отдельных классов.
\end{description}

CICIDS2017 в настоящей работе используется в режиме cross-evaluation для
проверки переносимости методики (см. \secref{subsec:unsw-vs-cic}).

\subsection{CSE-CIC-IDS2018}

CSE-CIC-IDS2018~\cite{sharafaldin2018toward} --- расширенная версия
CICIDS2017, размещённая в облачной инфраструктуре AWS, с увеличенным числом
пользователей-эмуляторов (50) и расширенным набором атак, включая botnet и
infiltration. Объём датасета на порядок больше CICIDS2017.

\begin{description}[font=\normalfont\itshape, leftmargin=1.5em]
    \item[Сильные стороны:] высокий объём данных, удобный для DL;
    более актуальные классы атак.
    \item[Ограничения:] унаследованы документированные дефекты CIC-семейства;
    высокая стоимость воспроизведения экспериментов из-за объёма.
\end{description}

\subsection{TON\_IoT}

TON\_IoT~\cite{moustafa2021toniot} ---  датасет той же исследовательской
группы UNSW Canberra Cyber, ориентированный на IoT и edge-инфраструктуры.
Содержит сетевой трафик, журналы Windows/Linux, телеметрию IoT-устройств.

\begin{description}[font=\normalfont\itshape, leftmargin=1.5em]
    \item[Сильные стороны:] современная (2019--2021) модель угроз;
    мультиисточниковый формат; преемственность с UNSW-NB15.
    \item[Ограничения:] узкая ориентация на IoT-сегмент, что ограничивает
    применимость к классической корпоративной LAN.
\end{description}

TON\_IoT в работе рассматривается как кандидат на расширение методики и
обсуждается в направлениях дальнейшего развития (раздел 4 главы 4 итоговой ВКР).

\subsection{Bot-IoT}

Bot-IoT~\cite{koroniotis2019botiot} --- датасет той же группы, акцентирующий
поведение ботнетов в IoT-сегменте. Содержит более 73 миллионов записей с
явным доминированием атакующего трафика.

\begin{description}[font=\normalfont\itshape, leftmargin=1.5em]
    \item[Сильные стороны:] большой объём; современная DDoS-таксономия.
    \item[Ограничения:] чрезмерная доля атакующего трафика (\(>97\%\)) делает
    датасет нерепрезентативным с точки зрения нормального профиля и требует
    специальной балансировки; ориентация на IoT.
\end{description}

\subsection{Kyoto 2006+}

Kyoto~\cite{song2011kyoto} --- датасет, собранный с продуктивных honeypot-узлов
сети Kyoto University. Содержит подмножество признаков, заимствованных из
KDDCup99, дополненное метками, основанными на реальных IDS-наблюдениях.

\begin{description}[font=\normalfont\itshape, leftmargin=1.5em]
    \item[Сильные стороны:] реальный (а не синтетический) трафик; длительный
    временной горизонт.
    \item[Ограничения:] узкое признаковое пространство, ограничивающее
    применимость DL-моделей; смещение в сторону сценариев honeypot, не
    отражающих корпоративную LAN.
\end{description}

\subsection{DARPA IDS Dataset}

DARPA IDS 1998/1999~\cite{lippmann2000darpa} --- исторический корпус,
послуживший основой для KDDCup99.

\begin{description}[font=\normalfont\itshape, leftmargin=1.5em]
    \item[Сильные стороны:] историческая значимость; первая публичная
    \emph{labeled}-выборка.
    \item[Ограничения:] устаревшая модель угроз (1998 г.); искусственный
    нормальный трафик; существенные методологические нарекания, изложенные
    в~\cite{sommer2010outside, moustafa2016eval}.
\end{description}

DARPA в качестве обучающего датасета сегодня неприемлем; в работе он
упоминается исключительно для обзорной полноты.

\subsection{LITNET-2020}

LITNET-2020~\cite{damasevicius2020litnet} --- датасет, собранный в реальной
магистральной сети Литвы (LITNET), содержащий аннотированные сетевые
flow-записи за 10 месяцев.

\begin{description}[font=\normalfont\itshape, leftmargin=1.5em]
    \item[Сильные стороны:] реальный продуктивный трафик; продолжительный
    временной горизонт.
    \item[Ограничения:] таксономия и объём отдельных классов несбалансированы;
    относительно слабое распространение в литературе; на момент написания
    работы --- ограниченное число публикаций с воспроизводимыми сравнениями.
\end{description}

\subsection{CICDDoS2019 (контекстная ссылка)}

CICDDoS2019~\cite{sharafaldin2019ddos} --- специализированный датасет,
посвящённый DDoS-атакам различных типов (UDP, SYN, MSSQL, NetBIOS, NTP, SNMP,
DNS-amp). Используется в работе как косвенная справочная таксономия при
формировании библиотеки атакующих сценариев для ИИ-злоумышленника.

\subsection{Сводное сравнение датасетов}

В \tabref{tab:datasets-compare} приведена сводная экспертная оценка
рассмотренных датасетов по критериям K1--K8. Оценки выставлены качественно по
шкале (\(-\), \(\circ\), \(+\)), где \(+\) обозначает соответствие критерию,
\(\circ\) --- частичное соответствие, \(-\) --- несоответствие. Источниками
оценок служат описанные выше публикации.

\begin{table}[H]
\centering
\caption{Сводное сравнение публично доступных датасетов NIDS}
\label{tab:datasets-compare}
\renewcommand{\arraystretch}{1.2}
\setlength{\tabcolsep}{4pt}
\small
\begin{tabularx}{\linewidth}{@{}lcccccccX@{}}
\toprule
\textbf{Датасет} & \textbf{K1} & \textbf{K2} & \textbf{K3} & \textbf{K4} &
\textbf{K5} & \textbf{K6} & \textbf{K7} & \textbf{K8: ключевые ограничения} \\
\midrule
KDDCup99~\cite{tavallaee2009nslkdd}            & $-$    & $-$ & $-$ & $\circ$ & $\circ$ & $-$    & $+$ & устаревшая модель угроз; дубликаты \\
NSL-KDD~\cite{tavallaee2009nslkdd}             & $-$    & $-$ & $\circ$ & $+$ & $\circ$ & $-$    & $+$ & таксономия 1998 г.; малый объём для DL \\
UNSW-NB15~\cite{moustafa2015unsw}              & $+$    & $+$ & $+$ & $\circ$ & $+$    & $+$    & $+$ & доминирование Generic; ct\_*-агрегаты \\
CICIDS2017~\cite{sharafaldin2018toward}        & $+$    & $+$ & $\circ$ & $\circ$ & $+$    & $+$    & $+$ & дефекты разметки; дубликаты \\
CSE-CIC-IDS2018~\cite{sharafaldin2018toward}   & $+$    & $+$ & $\circ$ & $\circ$ & $+$    & $+$    & $\circ$ & унаследованные дефекты; объём \\
TON\_IoT~\cite{moustafa2021toniot}             & $+$    & $\circ$ & $+$ & $\circ$ & $+$    & $+$    & $+$ & ориентация на IoT \\
Bot-IoT~\cite{koroniotis2019botiot}            & $\circ$& $-$ & $+$ & $-$    & $+$    & $+$    & $+$ & доля атак $>97\%$ \\
Kyoto 2006+~\cite{song2011kyoto}               & $\circ$& $\circ$ & $\circ$ & $\circ$ & $\circ$ & $\circ$ & $\circ$ & honeypot-смещение; узкие признаки \\
DARPA~\cite{lippmann2000darpa}                 & $-$    & $-$ & $\circ$ & $\circ$ & $-$    & $\circ$ & $+$ & 1998 г.; искусственный нормальный трафик \\
LITNET-2020~\cite{damasevicius2020litnet}      & $+$    & $+$ & $\circ$ & $\circ$ & $+$    & $+$    & $\circ$ & таксономический дисбаланс \\
\bottomrule
\end{tabularx}
\end{table}

Из приведённой таблицы видно, что наибольшее покрытие критериев у двух
датасетов: UNSW-NB15 и CICIDS2017. UNSW-NB15 превосходит CICIDS2017 по
качеству разметки и преемственности таксономии (последняя продолжена в
TON\_IoT и Bot-IoT того же научного коллектива), а также по более ранней
доступности и большей укоренённости в обзорной
литературе~\cite{khraisat2019survey, ahmad2021nids, ferrag2020dlcyber}.

\subsection{Обоснование выбора UNSW-NB15 в качестве основного датасета}
\label{subsec:unsw-justification}

С учётом сводного сравнения и характера задачи диссертационной работы UNSW-NB15
выбран в качестве основного датасета по следующим причинам.

\begin{enumerate}
    \item \textbf{Соответствие модели угроз.} Девять семейств атак UNSW-NB15
    покрывают рассмотренные в \tabref{tab:traffic-anomalies} типы
    несанкционированных действий (DoS, Reconnaissance, Exploits, Backdoors,
    Worms), а класс \emph{Generic} даёт фоновую массу атак, удобную для
    оценки FP-rate в условиях дисбаланса.
    \item \textbf{Совместимость с IPFIX/NetFlow.} Признаки UNSW-NB15
    структурно эквивалентны полям IPFIX-записей, что обеспечивает прямую
    переносимость методики на инфраструктуры реальных корпоративных сетей.
    \item \textbf{Стабильность разделений.} Использование официальных
    разделений \emph{training-set/testing-set} обеспечивает воспроизводимость
    результатов и сопоставимость с литературой.
    \item \textbf{Преемственность таксономии.} Близость UNSW-NB15 к TON\_IoT и
    Bot-IoT позволит в направлениях дальнейшего развития систематически
    расширять методику на edge- и IoT-сегменты без перестройки label-схемы.
    \item \textbf{Достаточный объём для DL.} Объём примерно \(2{,}54\,\text{млн}\)
    записей удовлетворяет требованию K5 и согласуется с минимальными порогами
    объёмов, рекомендуемыми для обучения CNN--LSTM-моделей в DL-NIDS-обзорах.
\end{enumerate}

\subsection{UNSW-NB15 vs CICIDS2017: схема cross-evaluation}
\label{subsec:unsw-vs-cic}

Для проверки переносимости методики (требование, вытекающее из ограничения
K8 и из обзора~\cite{pendlebury2019tesseract}) принят следующий протокол.

\begin{itemize}
    \item Модель обучается на UNSW-NB15 (training-set).
    \item Производится унификация подмножества признаков, общих для UNSW-NB15
    и CICIDS2017 (\emph{duration}, \emph{bytes/pkts in/out}, \emph{flags},
    \emph{protocol}, \emph{service}); список общих признаков
    фиксируется в главе~2 настоящей работы.
    \item Тест проводится на размеченной выборке CICIDS2017 без дообучения и
    с дообучением на 10\% выборки CICIDS2017 (контрольный режим transfer
    learning).
    \item Оценивается падение F1, PR-AUC и FP-rate. Падение интерпретируется
    как мера переносимости (\eng{generalization gap}).
\end{itemize}

\subsection{Угрозы валидности и пути их минимизации}

С учётом замечаний~\cite{sommer2010outside, pendlebury2019tesseract,
gharib2016evalframework} в работе предусмотрены следующие меры по снижению
угроз внутренней и внешней валидности экспериментов:

\begin{enumerate}
    \item Использование официального train/test-разделения UNSW-NB15 без
    дополнительного перемешивания между разделами.
    \item Контроль временного смещения (\emph{temporal split}) при
    формировании скользящих окон --- запрет утечек ``из будущего''.
    \item Прогон экспериментов с фиксированными seed для воспроизводимости
    (\(\geq 5\) seed на эксперимент).
    \item Кросс-проверка результатов на CICIDS2017 как независимом датасете.
    \item Учёт стоимости ложных срабатываний на единицу времени
    (FP-per-time), что компенсирует потенциально завышенную ROC-AUC при
    дисбалансе.
\end{enumerate}

\subsection{Промежуточные выводы по разделу}

\begin{enumerate}
    \item Анализ десяти публично доступных корпусов трафика по восьми
    критериям K1--K8 показал, что наибольшее покрытие критериев имеют
    UNSW-NB15 и CICIDS2017; остальные либо устарели по модели угроз
    (KDD-семейство, DARPA), либо ориентированы на IoT-сегмент
    (TON\_IoT, Bot-IoT), либо имеют известные методологические дефекты
    разметки.
    \item UNSW-NB15 превосходит CICIDS2017 по качеству разметки,
    преемственности таксономии и совместимости со стандартом IPFIX,
    что обосновывает его выбор в качестве основного обучающего и тестирующего
    корпуса.
    \item CICIDS2017 принят как вспомогательный корпус для cross-evaluation;
    разработана схема унификации общих признаков и протокол оценки переносимости
    методики.
    \item Зафиксирован набор мер по минимизации угроз валидности
    (фиксированные разделения, временной split, повторение экспериментов
    по seed, кросс-проверка), что соответствует современным требованиям к
    обучающим протоколам в области ML-NIDS.
\end{enumerate}

% =====================================================================
%  Stage 1, раздел 3.
%  Обзор подходов к моделированию атак, генеративных моделей
%  и обоснование архитектуры ИИ-злоумышленника.
% =====================================================================

\section{Обзор подходов к моделированию атак и обоснование архитектуры ИИ-злоумышленника}

\subsection{Назначение ИИ-злоумышленника в системе оценки NIDS}

Под ИИ-злоумышленником (далее --- агентом) в работе понимается активный
программный компонент, формирующий поток сетевых событий с управляемым
переключением режимов \emph{норма / атака / дрейф} и подающий этот поток на
вход детектора. Назначение агента --- обеспечить воспроизводимое и
параметризованное окружение для оценки качества и устойчивости детектора в
условиях, приближенных к эксплуатационным.

В отличие от пассивных датасетов, агент позволяет:
\begin{itemize}
    \item получить контролируемую \emph{ground truth} в потоковом сценарии,
    включая моменты переключения режимов;
    \item варьировать частоты атак, длительность сессий и параметры сетевых
    шаблонов в широких пределах;
    \item моделировать дрейф распределений, проверяя устойчивость детектора
    и работу drift-мониторинга;
    \item обеспечить независимость тестового сценария от датасета,
    использованного для обучения, что снижает риск переобучения на специфике
    UNSW-NB15.
\end{itemize}

Отдельно подчеркнём, что речь идёт о \emph{симуляционном агенте на уровне
flow-признаков}: целью не является эмуляция реальных атак на работающую
инфраструктуру. Это разделение существенно с точки зрения этики и
юридического статуса исследования.

\subsection{Подходы к моделированию атак: систематика}

В литературе по ML-NIDS выделяются три устойчивых семейства подходов к
формированию атакующего трафика для тестирования детекторов.

\begin{enumerate}
    \item \textbf{Шаблонная (\eng{rule-based, template-based}) генерация.}
    Атакующие сценарии задаются параметрическими шаблонами: распределения
    межпакетных интервалов, длины потоков, наборы флагов TCP, скорости
    передачи. Применение --- широкое: от классических работ DARPA/KDD до
    инструментов \eng{traffic replay} вроде Tcpreplay и
    Hping3~\cite{lippmann2000darpa, gharib2016evalframework}.
    \item \textbf{Генеративные модели.} Использование GAN, cGAN, WGAN,
    VAE и автоэнкодеров для синтеза реалистичных flow-векторов и
    последовательностей. Эталонная работа --- Ring и
    соавт.~\cite{ring2019flowgan}; вариационные подходы для IoT-сетей
    предложены в~\cite{lopez2017cvae}.
    \item \textbf{Adversarial-агенты на основе обучения с подкреплением.}
    Агент учится формировать последовательности действий, обходящих
    детектор, в форме, аналогичной white-/grey-box атаке на классификатор;
    эталонная работа --- Apruzzese и соавт.~\cite{apruzzese2020adversarial}.
\end{enumerate}

В \tabref{tab:attacker-approaches} приведено сводное теоретическое сравнение
этих семейств по пяти эксплуатационно значимым критериям.

\begin{table}[H]
\centering
\caption{Сравнение подходов к моделированию атак для оценки NIDS}
\label{tab:attacker-approaches}
\renewcommand{\arraystretch}{1.2}
\begin{tabularx}{\linewidth}{@{}lXXX@{}}
\toprule
\textbf{Критерий} & \textbf{Шаблонная} & \textbf{Генеративная (GAN/VAE)} &
\textbf{Adversarial-RL} \\
\midrule
Воспроизводимость       & высокая  & средняя; зависит от стабильности обучения &
низкая; зависит от инициализации политики \\
\addlinespace
Семантическая
интерпретируемость      & прямая (параметры~$\to$~атака) & косвенная (через
обученное распределение) & низкая (политика — «чёрный ящик») \\
\addlinespace
Вариативность           & ограничена шаблонами; усиливается рандомизацией параметров &
высокая; ограничена режимом коллапса GAN & высокая, но смещена в сторону
обхода конкретного детектора \\
\addlinespace
Стоимость разработки    & низкая & высокая (стабильное обучение GAN) &
очень высокая (среда + reward design + обучение) \\
\addlinespace
Совместимость с этикой
и охватом ВКР           & полная & требует валидации реалистичности & требует
явного контроля; возможна оптимизация под обход детектора \\
\bottomrule
\end{tabularx}
\end{table}

\subsection{Шаблонная генерация: математическая модель}

Шаблон атаки описывается параметрическим вектором
\(\theta \in \Theta\), задающим распределения базовых характеристик
flow-сессий: длительность \(d \sim F_d(\cdot;\theta)\), скорость пакетов
\(\lambda \sim F_\lambda(\cdot;\theta)\), доля флагов SYN/ACK и др. Поток
наблюдений на интервале \([0, T]\) задаётся точечным процессом
\(\mathcal{N}(t;\theta)\), параметры которого определяются типом атаки.
Например, для DDoS volumetric SYN-flood:
\begin{equation}
\lambda(t) = \lambda_0 \cdot \mathbb{I}[t \in [t_a, t_a + \Delta]] + \lambda_n,
\label{eq:ddos-rate}
\end{equation}
где \(\lambda_n\) --- интенсивность фонового потока, \(\lambda_0\) ---
интенсивность атаки, \([t_a, t_a + \Delta]\) --- временной интервал атаки.
Шаблон позволяет изменять \(\lambda_0\), \(\Delta\) и долю \emph{SYN}-флагов,
получая семейство сценариев одной и той же атаки.

Аналогичные шаблоны определены для остальных классов: Reconnaissance задаётся
заметанием destination-портов с фиксированной интенсивностью обращений к
одному источнику; Brute Force --- регулярными подключениями к одному порту с
малым числом байт за сессию; Exfiltration --- длинной устойчивой передачей с
дисбалансом \(b_{out} \gg b_{in}\); Internal Reconnaissance --- широким
фан-аутом потоков от одного источника к множеству destinations с типичной
поведенческой подписью.

Преимущество шаблонного подхода --- прозрачность ground truth: момент атаки,
её тип и параметры однозначно зафиксированы конфигурацией. Ограничение ---
ограниченная вариативность; компенсируется введением рандомизации параметров.

\subsection{Генеративные модели для синтеза трафика: теоретическая база}

\subsubsection{GAN: формализм}

Классическая постановка GAN~\cite{goodfellow2014gan}:
\begin{equation}
\min_{G}\max_{D}
V(D,G) = \mathbb{E}_{x \sim p_{\text{data}}}[\log D(x)] +
         \mathbb{E}_{z \sim p_z}[\log(1 - D(G(z)))],
\label{eq:gan}
\end{equation}
где \(G\colon \mathcal{Z} \to \mathcal{X}\) --- генератор,
\(D\colon \mathcal{X} \to [0,1]\) --- дискриминатор,
\(p_z\) --- априорное распределение в латентном пространстве.
Минимакс \eqnref{eq:gan} в идеале сводится к равенству распределений
\(G(z) \sim p_{\text{data}}\) при оптимальном дискриминаторе.

\subsubsection{Conditional GAN}

Conditional GAN~\cite{mirza2014cgan} вводит обусловливающую переменную \(c\)
(тип атаки):
\begin{equation}
\min_{G}\max_{D}
V(D,G) = \mathbb{E}_{x,c}[\log D(x,c)] +
         \mathbb{E}_{z,c}[\log(1 - D(G(z,c), c))].
\label{eq:cgan}
\end{equation}
Это позволяет генерировать целенаправленные атакующие сценарии конкретного
класса, что и составляет содержание подхода
Ring и соавт.~\cite{ring2019flowgan}.

\subsubsection{WGAN-GP}

Обучение GAN страдает от нестабильности и режима коллапса (\eng{mode
collapse}). Wasserstein GAN~\cite{arjovsky2017wgan} заменяет JS-расстояние
расстоянием Earth Mover (Wasserstein-1):
\begin{equation}
W(p, q) = \inf_{\gamma \in \Pi(p, q)} \mathbb{E}_{(x,y) \sim \gamma}\|x - y\|.
\label{eq:wgan}
\end{equation}
Расширение WGAN-GP~\cite{gulrajani2017wgangp} вводит градиентный штраф,
обеспечивающий 1-липшицевость дискриминатора:
\begin{equation}
\mathcal{L}_{\text{GP}} = \lambda\,
\mathbb{E}_{\hat{x}}\bigl[(\|\nabla_{\hat{x}} D(\hat{x})\|_2 - 1)^2\bigr].
\label{eq:wgangp}
\end{equation}
WGAN-GP --- стандартный выбор для синтеза табличных flow-данных.

\subsubsection{Применение GAN-семейства в NIDS}

Систематический обзор Ring и соавт.~\cite{ring2019flowgan} показал, что
GAN способен синтезировать реалистичные flow-векторы с распределениями,
неотличимыми по ряду статистических критериев. Однако на практике
зафиксированы:

\begin{itemize}
    \item Чувствительность качества к гиперпараметрам и инициализации.
    \item Режим коллапса для редких классов (Worms, Backdoors), где
    обучающих данных мало.
    \item Сложность валидации семантической согласованности
    сгенерированного трафика (например, согласованности TCP-флагов и
    длительности).
\end{itemize}

Эти ограничения существенны в условиях ВКР, где требуется обеспечить
воспроизводимое тестирование и однозначную ground truth.

\subsection{Adversarial-агенты на основе обучения с подкреплением}

Apruzzese и соавт.~\cite{apruzzese2020adversarial} рассматривают атакующего
как агента в среде, где \(\text{state}\) --- текущее состояние
взаимодействия с детектором, \(\text{action}\) --- модификация параметров
исходящего трафика, \(\text{reward}\) --- успех обхода детектора.
Стандартная политика обучается алгоритмами PPO, DQN или DDPG.

С точки зрения настоящей работы такой подход:
\begin{itemize}
    \item даёт мощный инструмент для аудита устойчивости детектора;
    \item трудоёмок в реализации (требует среды, симулирующей ответ
    детектора в режиме обучения);
    \item смещает поведение агента в сторону обхода конкретного
    детектора, что не совпадает с задачей формирования \emph{внешних}
    тестовых сценариев;
    \item требует расширенного этического контроля и явного отказа от
    обобщения на реальные системы.
\end{itemize}

В рамках ВКР этот подход рассматривается как направление дальнейшего
развития, но не как реализуемый компонент.

\subsection{Concept drift: моделирование и оценка}

Феномен \emph{concept drift} --- изменение распределений данных и/или
условной зависимости меток от признаков во времени --- является основным
источником деградации NIDS в эксплуатации~\cite{gama2014drift,
pendlebury2019tesseract, andresini2021insomnia}.

Различают четыре типа дрейфа~\cite{gama2014drift}:
\begin{enumerate}
    \item \textbf{Sudden drift} --- резкое изменение распределения
    \(p(x, y)\);
    \item \textbf{Incremental drift} --- постепенное смещение во времени;
    \item \textbf{Gradual drift} --- частая смена режимов между двумя
    распределениями;
    \item \textbf{Recurring drift} --- циклически возвращающиеся режимы.
\end{enumerate}

Для целей настоящей работы дрейф моделируется как преобразование
\(\mathcal{T}\colon X \to X\), применяемое к подмножеству признаков. Три
варианта инжектора реализуются в Этапе 3 и используются в Эксперименте 5
Этапа 4:

\begin{itemize}
    \item \textbf{Covariate shift}:
    \(x' = \mathcal{T}_{\text{cov}}(x; \mu, \sigma) = \mu + \sigma \odot x\),
    что соответствует изменению масштаба и положения признаков без
    изменения метки;
    \item \textbf{Prior shift}: изменение долей классов в потоке
    через политику переключения режимов \emph{ATTACK(type)} в
    конечном автомате агента;
    \item \textbf{Concept drift}: изменение зависимости \(p(y \mid x)\)
    через модификацию параметров шаблонов атак (например, понижение
    \(\lambda_0\) DDoS до уровней, граничащих с нормой).
\end{itemize}

Контроль интенсивности дрейфа осуществляется через индексы
\(\PSI\)~\eqnref{eq:psi} и \(\MMD\)~\eqnref{eq:mmd}, что обеспечивает
сопоставимость экспериментов.

\subsection{Архитектура ИИ-злоумышленника, принимаемая в работе}

С учётом проведённого сравнения и условий, формируемых задачей ВКР, в работе
принимается \textbf{гибридная} архитектура агента:

\[
\boxed{\;\text{Templates} \;+\; \text{FSM-планировщик} \;+\; \text{Drift-инжектор}\;}
\]

\noindent без использования генеративных моделей в качестве реализуемого
компонента. Генеративные подходы (GAN/cGAN/WGAN-GP) рассматриваются как
теоретическая альтернатива, обоснованно отвергаемая по следующим причинам.

\begin{enumerate}
    \item \textbf{Воспроизводимость.} Шаблонный агент даёт однозначную
    \emph{ground truth} (момент, тип, параметры атаки). GAN-агент, напротив,
    требует фиксации seeds, состояния обучения и контроля режима коллапса,
    что существенно усложняет научную верификацию результатов.
    \item \textbf{Стабильность обучения.} В условиях ВКР с ограниченным числом
    атакующих образцов (особенно для редких классов в UNSW-NB15) GAN-обучение
    подвержено режиму коллапса, что приводит к нерепрезентативным сценариям и
    систематическим ошибкам экспериментальной части.
    \item \textbf{Семантическая согласованность.} Шаблоны явно поддерживают
    согласованность между признаками (например, TCP-флагами и
    длительностью), тогда как GAN-генерация на табличных flow-данных требует
    отдельных верификационных проверок (TSTR, маргиналы, корреляции).
    \item \textbf{Привязка к таксономии.} Шаблоны позволяют детерминированно
    соответствовать классам UNSW-NB15, что упрощает интерпретацию ошибок
    детектора и анализ FP/FN.
\end{enumerate}

При этом сохраняемое в работе теоретическое описание GAN-семейства
обеспечивает полноту обзора и явно фиксирует направление дальнейшего
развития методики --- расширение шаблонной библиотеки за счёт обученной
cGAN-составляющей.

\subsection{Конечно-автоматная политика агента}

Состояния планировщика образуют конечный автомат:
\[
\mathcal{S} = \{\text{NORMAL}\} \cup
\{\text{ATTACK}(c)\}_{c \in \mathcal{C}} \cup
\{\text{DRIFT}(k)\}_{k \in \mathcal{K}},
\]
где \(\mathcal{C}\) --- множество типов атак из таксономии UNSW-NB15
(в реализации --- подмножество с реалистичной моделью flow-проявлений),
\(\mathcal{K}\) --- множество типов дрейфа (covariate / prior / concept).

Переходы между состояниями определяются дискретной политикой
\(\pi\colon \mathcal{S} \to \Delta(\mathcal{S})\), параметризованной
\emph{вероятностями переключения} и \emph{минимальными длительностями
сессий} для каждого состояния. Такая параметризация обеспечивает:
\begin{itemize}
    \item возможность задавать долю атакующего трафика в потоке (\emph{attack
    ratio});
    \item возможность задавать частоту смены режимов (\emph{flicker rate});
    \item полную воспроизводимость прогона при фиксированном seed.
\end{itemize}

Формально политика для перехода из NORMAL имеет вид
\[
\pi(\text{NORMAL}) =
\bigl((1-p_a-p_d)\!\cdot\!\text{NORMAL}\bigr)
\oplus \sum_{c}\!p_a^{(c)}\!\cdot\!\text{ATTACK}(c)
\oplus \sum_{k}\!p_d^{(k)}\!\cdot\!\text{DRIFT}(k),
\]
где \(\oplus\) --- смесь распределений. Минимальные длительности
\(d_{\min}^{s}\) для каждого состояния гарантируют отсутствие
``мерцающих'' переключений, нерепрезентативных в эксплуатации.

\subsection{Ground truth, метаданные и протокол валидации агента}

Каждое сгенерированное окно сопровождается записью \(\langle t, s, \theta\rangle\),
где \(t\) --- метка времени, \(s\) --- состояние FSM, \(\theta\) ---
параметры шаблона. Этот журнал служит неизменной \emph{ground truth} для
вычисления метрик (особенно time-to-detect и FP-per-time).

Протокол валидации агента включает:
\begin{enumerate}
    \item \textbf{Маргинальное соответствие.} Сравнение распределений ключевых
    признаков (\emph{duration, sbytes, dbytes, sttl, dur}) в режиме NORMAL
    с распределениями нормального трафика UNSW-NB15 по тестам
    Колмогорова--Смирнова и MMD. Цель --- подтвердить, что фон агента не
    тривиально отличается от референсной нормы.
    \item \textbf{Семантическую согласованность атак.} Проверка соответствия
    флагов TCP, длительности и интенсивности шаблонных атак данным UNSW-NB15
    для соответствующего класса.
    \item \textbf{Контроль интенсивности дрейфа.} Расчёт PSI/MMD между
    эталонным и текущим окном, попадание в заданные диапазоны при
    включении DRIFT-состояний.
\end{enumerate}

\subsection{Этические и юридические аспекты}

Все процедуры моделирования атак выполняются исключительно в составе
автономного программного стенда; целью не является эмуляция реальных атак
на работающие сети или сервисы. Сценарии формируются в виде синтетических
flow-векторов, не передаваемых в реальную сетевую среду. В работе
отсутствуют операции, направленные на обход эксплуатируемых средств защиты
третьих лиц, и отсутствуют рекомендации по практическому применению
атакующих техник. Это соответствует общепринятой в академическом сообществе
практике исследований adversarial-NIDS~\cite{apruzzese2020adversarial,
ring2019flowgan}.

\subsection{Промежуточные выводы по разделу}

\begin{enumerate}
    \item Систематизированы три семейства подходов к моделированию атак
    для оценки NIDS: шаблонные, генеративные и adversarial-агенты на
    обучении с подкреплением. Каждое семейство охарактеризовано по пяти
    критериям; отмечено, что шаблонный подход обеспечивает наивысшую
    воспроизводимость и интерпретируемость, но ограничен по вариативности.
    \item Описана математическая база генеративного семейства (GAN, cGAN,
    WGAN-GP), приведены целевые функции и зафиксированы ограничения
    применимости в условиях ВКР: нестабильность обучения, режим коллапса
    для редких классов, сложность валидации семантической согласованности.
    \item Обоснован выбор гибридной архитектуры агента: \emph{шаблоны +
    конечный автомат + drift-инжектор} без cGAN. Описаны математические
    модели шаблонов основных классов атак (DDoS, Reconnaissance, Brute Force,
    Exfiltration, Internal Reconnaissance) и формальная политика FSM.
    \item Зафиксированы три типа моделируемого дрейфа (covariate, prior,
    concept) и метрики его контроля (PSI, MMD), что обеспечивает связность
    с протоколом экспериментов главы~4.
    \item Сформулирован протокол валидации агента (маргинальное соответствие,
    семантическая согласованность, контроль интенсивности дрейфа) и
    зафиксированы этические рамки исследования.
\end{enumerate}

% =====================================================================
%  Stage 1, раздел 4.
%  Формализация требований к системе и гипотезы исследования.
% =====================================================================

\section{Требования к системе и формализация гипотезы исследования}

\subsection{Концептуальная модель системы}

В обобщённой постановке разрабатываемая система состоит из трёх подсистем,
взаимодействие которых определено протоколом обмена событиями.

\begin{enumerate}
    \item \textbf{Подсистема предобработки и формирования окон.}
    Принимает поток сетевых соединений в формате IPFIX-совместимых
    flow-записей, выполняет валидацию схемы, очистку, нормализацию и
    формирование скользящих окон фиксированной длины.
    \item \textbf{Подсистема детекции (нейросетевая модель).}
    Принимает окна, формирует скоринг \(s(x)\) и --- после применения
    откалиброванного порога \(\tau\) --- бинарное решение об аномалии,
    сопровождаемое метаданными для приоритизации алертов.
    \item \textbf{Подсистема активного тестирования (ИИ-злоумышленник).}
    Воспроизводит контролируемый поток событий с переключением режимов
    \emph{норма / атака / дрейф} и подаёт его в подсистему предобработки.
\end{enumerate}

В составе SOC рассматриваемая система занимает уровень NTA-аналитики и
взаимодействует с источниками сетевой телеметрии и подсистемой алертирования
(SIEM/SOAR), что согласуется с \secref{subsec:soc-integration} раздела~1.

\subsection{Функциональные требования (\textbf{FR})}

\begin{description}[font=\normalfont\bfseries, leftmargin=2.5em]
    \item[FR-1.] Система должна принимать на вход flow-записи, совместимые
    со схемой UNSW-NB15 (49 признаков) и подмножеством, общим с CICIDS2017
    (для cross-evaluation), с явной валидацией типов и диапазонов.
    \item[FR-2.] Система должна выполнять предобработку признаков:
    очистку пропусков, log-преобразование признаков с положительной
    асимметрией, robust-scaling, кодирование категориальных признаков,
    формирование скользящих окон.
    \item[FR-3.] Система должна обеспечивать обучение, валидацию и сохранение
    весов нейросетевой модели CNN--LSTM (основной) и моделей-альтернатив
    (1D-CNN, LSTM, BiLSTM, GRU, TCN, Transformer-encoder, Autoencoder, VAE,
    MLP) и базовых моделей (Logistic Regression, Random Forest, XGBoost,
    Isolation Forest, One-Class SVM).
    \item[FR-4.] Система должна формировать скоринг \(s(x)\) и применять
    откалиброванный порог \(\tau\), полученный по целевому уровню FP-rate
    или FP-per-time.
    \item[FR-5.] Система должна формировать алерты с метаданными:
    идентификатор окна, метка времени, скоринг, предсказанный класс,
    уровень приоритета, признаковый профиль для аналитика.
    \item[FR-6.] Система должна предоставлять REST-интерфейс на базе FastAPI
    для подачи окна (или последовательности окон) и получения скоринга,
    решения и метаданных.
    \item[FR-7.] Система должна предоставлять минимальный пользовательский
    интерфейс на базе Streamlit для отображения потока алертов, метрик
    качества и состояния drift-мониторинга.
    \item[FR-8.] Система должна включать ИИ-злоумышленника на основе
    \emph{шаблонов + конечного автомата + drift-инжектора} и обеспечивать
    конфигурируемые сценарии тестирования.
    \item[FR-9.] Система должна реализовывать drift-мониторинг по индексам
    PSI, KL и MMD относительно эталонного окна нормы.
    \item[FR-10.] Все эксперименты должны фиксировать seed
    (\(\geq 5\) независимых запусков) и записывать метрики и
    конфигурации в журналы (CSV/JSON) для воспроизведения.
\end{description}

\subsection{Нефункциональные требования (\textbf{NFR})}

\begin{description}[font=\normalfont\bfseries, leftmargin=2.5em]
    \item[NFR-1.] \emph{Качество детекции на UNSW-NB15.}
    F1 \(\geq 0{,}90\) и PR-AUC \(\geq 0{,}95\) на официальном
    тестовом разделении UNSW-NB15 при бинарной постановке.
    \item[NFR-2.] \emph{Уровень ложных срабатываний.} FPR \(\leq 1\%\)
    при Recall \(\geq 0{,}85\) после калибровки порога; контроль
    FP-per-time для оценки эксплуатационной нагрузки.
    \item[NFR-3.] \emph{Эксплуатационная задержка.} Inference latency
    \(\leq 50\) мс на одно окно при batch-size = 1 на CPU-конфигурации
    среднего класса.
    \item[NFR-4.] \emph{Пропускная способность.} Throughput
    \(\geq 1000\) flow-записей в секунду (EPS) при батчировании
    окон на CPU-конфигурации среднего класса.
    \item[NFR-5.] \emph{Time-to-detect.} TTD \(\leq 1\) окна с момента
    появления атаки в потоке (для шаблонов с явными flow-проявлениями
    --- DDoS, Port Scan, Brute Force).
    \item[NFR-6.] \emph{Устойчивость к дрейфу.} Снижение F1 при включении
    drift-инжектора с интенсивностью PSI~\(\leq 0{,}25\) не должно
    превышать абсолютной величины \(\Delta\,\text{F1} = 0{,}10\) до
    срабатывания механизма адаптации/повторной калибровки.
    \item[NFR-7.] \emph{Перенос на CICIDS2017.} В режиме \emph{train on
    UNSW, test on CICIDS2017} (без дообучения) \(\Delta\,\text{F1}\)
    относительно качества на UNSW \(\leq 0{,}20\); с дообучением на 10\%
    выборки CICIDS2017 \(\Delta\,\text{F1} \leq 0{,}10\).
    \item[NFR-8.] \emph{Воспроизводимость.} Полный pipeline (от raw-CSV
    до итоговых метрик и графиков отчёта) запускается одной командой
    Makefile при фиксированных seed.
    \item[NFR-9.] \emph{Стабильность.} Дисперсия F1 по seed
    (\(n \geq 5\)) удовлетворяет \(\sigma_{F1} \leq 0{,}01\).
\end{description}

\subsection{Метрики и протокол оценки}

В соответствии с \secref{sec:metrics} раздела~1, для оценки разработанной
методики применяются три группы метрик.

\begin{enumerate}
    \item \textbf{Классификационные:} Precision, Recall, F1, MCC, FPR.
    \item \textbf{Интегральные:} PR-AUC (приоритетно), ROC-AUC.
    \item \textbf{Эксплуатационные:} time-to-detect, latency, throughput,
    FP-per-time.
\end{enumerate}

Дополнительно фиксируются:
\begin{itemize}
    \item Калибровка предсказаний (ECE);
    \item Метрики дрейфа (PSI, MMD относительно эталонного окна);
    \item Стабильность по seed (\(\sigma_{F1}\), \(\sigma_{\text{PR-AUC}}\)).
\end{itemize}

Протокол сравнения моделей: одинаковое разделение, одинаковый препроцессинг,
одинаковый порог калибровки (по целевому FP-rate). Статистическая значимость
различий между моделями оценивается парным критерием Уилкоксона по
результатам \(n \geq 5\) запусков.

\subsection{Допущения и ограничения исследования}

\begin{enumerate}
    \item \textbf{Допущение об отсутствии deep packet inspection.} Анализ
    выполняется на flow-уровне; полезная нагрузка пакетов не учитывается.
    Такое допущение согласуется с реальной практикой применения NIDS в
    шифрованном трафике и со стандартом IPFIX~\cite{rfc7011}.
    \item \textbf{Допущение об автономности активного тестирования.}
    ИИ-злоумышленник работает в рамках программного стенда и не порождает
    реальный сетевой трафик за пределами стенда.
    \item \textbf{Ограничение по числу классов.} Мультиклассовая постановка
    рассматривается дополнительно (см. главу~2), однако основная гипотеза
    проверяется в бинарной постановке \emph{normal/attack}.
    \item \textbf{Ограничение по типам дрейфа.} Моделируются три типа
    дрейфа (covariate, prior, concept). Recurring drift и более сложные
    нестационарные режимы отнесены к направлениям дальнейшего развития.
    \item \textbf{Ограничение по DL-фреймворку.} Модели реализуются на
    PyTorch, классические baseline --- на scikit-learn и XGBoost. Сравнение
    с другими фреймворками вне области исследования.
\end{enumerate}

\subsection{Формализация гипотезы исследования}

\subsubsection{Основная гипотеза}

\begin{statement}[Основная гипотеза, $H_0^{\text{main}}$]
Существует нейросетевая модель \(M^*\) с архитектурой \emph{CNN--LSTM},
обученная на flow-представлении датасета UNSW-NB15 в скользящих окнах,
такая, что на официальном тестовом разделении датасета:
\begin{enumerate}
    \item F1 \(\geq 0{,}90\) и PR-AUC \(\geq 0{,}95\),
    \item FPR \(\leq 1\%\) при Recall \(\geq 0{,}85\),
\end{enumerate}
и эти показатели статистически значимо превосходят (\(p < 0{,}05\) по парному
критерию Уилкоксона) показатели ансамбля моделей-baseline (Logistic
Regression, Random Forest, XGBoost, Isolation Forest, One-Class SVM, MLP)
при равных условиях обучения и оценки.
\end{statement}

\subsubsection{Дополнительная гипотеза о роли активного тестирования}

\begin{statement}[$H_0^{\text{aux}}$]
Применение ИИ-злоумышленника в форме \emph{шаблоны + конечный автомат +
drift-инжектор} к разработанному прототипу позволяет:
\begin{enumerate}
    \item обнаруживать аномалии в потоковом сценарии с TTD \(\leq 1\) окна
    для шаблонных атак с явными flow-проявлениями (DDoS, Port Scan, Brute
    Force);
    \item фиксировать снижение качества детекции при контролируемом
    дрейфе (PSI~\(\leq 0{,}25\)) в пределах \(\Delta\,\text{F1} \leq 0{,}10\);
    \item систематически идентифицировать классы со сниженной полнотой
    обнаружения и связывать ошибки FN с конкретными состояниями FSM, что
    обеспечивает дифференциальный анализ устойчивости детектора, недоступный
    при использовании только статического датасета.
\end{enumerate}
\end{statement}

\subsubsection{Гипотеза переносимости}

\begin{statement}[$H_0^{\text{transfer}}$]
Модель \(M^*\), обученная на UNSW-NB15, демонстрирует на CICIDS2017
\(\Delta\,\text{F1} \leq 0{,}20\) без дообучения и
\(\Delta\,\text{F1} \leq 0{,}10\) с дообучением на 10\% выборки CICIDS2017
относительно качества на UNSW-NB15.
\end{statement}

\subsection{Критерии подтверждения и опровержения гипотезы}

Гипотеза \(H_0^{\text{main}}\) считается подтверждённой, если на
\(n \geq 5\) запусках с независимыми seed выполнены условия NFR-1, NFR-2 и
получено статистически значимое превосходство над ансамблем baseline
(\(p < 0{,}05\)).

Гипотеза \(H_0^{\text{aux}}\) подтверждается при выполнении NFR-5 для
шаблонных атак, NFR-6 для drift-сценариев и при наличии в результатах
систематического разделения ошибок FN по состояниям FSM.

Гипотеза \(H_0^{\text{transfer}}\) подтверждается при выполнении NFR-7.

В случае невыполнения какого-либо из критериев гипотеза подвергается
уточнению по итогам анализа: возможны изменения архитектуры (переход к
BiLSTM~+~Attention, Transformer-encoder), параметров препроцессинга,
порога калибровки или политики FSM. Каждое такое изменение фиксируется в
журнале экспериментов.

\subsection{Промежуточные выводы по разделу}

\begin{enumerate}
    \item Сформирована концептуальная трёхуровневая модель системы
    (предобработка, детекция, активное тестирование) с явным определением
    интерфейсов и места в структуре SOC.
    \item Определены десять функциональных и девять нефункциональных
    требований с количественными порогами по качеству детекции, задержкам,
    устойчивости к дрейфу и переносимости. Все пороги связаны с метриками
    \secref{sec:metrics} и обоснованы с учётом эксплуатационной природы
    NIDS.
    \item Сформулирована основная гипотеза исследования \(H_0^{\text{main}}\)
    о превосходстве CNN--LSTM на UNSW-NB15 над ансамблем baseline, а также
    дополнительные гипотезы о роли ИИ-злоумышленника
    (\(H_0^{\text{aux}}\)) и о переносимости методики на CICIDS2017
    (\(H_0^{\text{transfer}}\)) с явными критериями подтверждения и
    опровержения.
    \item Определены допущения и ограничения исследования (flow-уровень
    анализа, программно-стендовое тестирование, бинарная постановка как
    основная), что задаёт область применимости разработанной методики и
    границы интерпретации экспериментальных результатов.
\end{enumerate}


% ----------- Заключение -----------
\section*{ЗАКЛЮЧЕНИЕ}
\addcontentsline{toc}{section}{ЗАКЛЮЧЕНИЕ}

В результате выполнения Этапа 1 получены следующие основные результаты.

\begin{enumerate}
    \item Выполнен систематический анализ современных подходов к
    обнаружению аномалий и атак в сетевом трафике. Установлено, что
    нейросетевые методы --- в особенности гибридные архитектуры
    CNN--LSTM --- являются наиболее перспективными для flow-уровневого
    анализа благодаря способности одновременно учитывать локальные
    паттерны и временные зависимости. Систематизированы три группы
    метрик оценки (классификационные, интегральные, эксплуатационные)
    с приоритетом PR-AUC и FP-per-time как наиболее адекватных для
    дисбалансированных NIDS-задач.
    \item Проведено сравнительное исследование десяти публично доступных
    датасетов NIDS по восьми формализованным критериям. Установлено, что
    UNSW-NB15 удовлетворяет критериям актуальности, репрезентативности,
    качества разметки, пригодности для DL и flow-анализа в наибольшей
    степени, и выбран в качестве основного. CICIDS2017 принят как
    вспомогательный для cross-evaluation; разработана схема унификации
    общих признаков и протокол оценки переносимости методики.
    \item Систематизированы подходы к моделированию атак (шаблонные,
    генеративные, adversarial-RL) и приведена их математическая база.
    На основании сопоставления по пяти эксплуатационным критериям
    (воспроизводимость, интерпретируемость, вариативность, стоимость
    разработки, этическая приемлемость) обоснован выбор гибридной
    архитектуры ИИ-злоумышленника: \emph{шаблоны + конечный автомат +
    drift-инжектор} без cGAN-составляющей. Сформулирован протокол
    валидации агента и зафиксированы этические рамки исследования.
    \item Сформирована концептуальная модель системы, включающая
    подсистемы предобработки, детекции и активного тестирования.
    Определены 10 функциональных и 9 нефункциональных требований с
    количественными порогами. Формализованы три гипотезы исследования:
    основная (превосходство CNN--LSTM над baseline на UNSW-NB15),
    дополнительная (роль ИИ-злоумышленника в оценке устойчивости) и
    гипотеза переносимости (на CICIDS2017). Для каждой гипотезы заданы
    критерии подтверждения и опровержения.
\end{enumerate}

Полученные результаты используются как входные данные Этапа 2:
теоретическое сравнение архитектур, формирование данных и протокола
экспериментов.

% Метки для подсчёта последних элементов
\newcounter{lastfigvalue}
\newcounter{lasttabvalue}
\newcounter{lastrefvalue}
\label{lastfig}
\label{lasttab}
\label{lastref}

% ----------- Список использованных источников -----------
\printbibliography[title={СПИСОК ИСПОЛЬЗОВАННЫХ ИСТОЧНИКОВ}]

% Метка последней страницы
\label{LastPage}

\end{document}
